% generated by GAPDoc2LaTeX from XML source (Frank Luebeck)
\documentclass[a4paper,11pt]{report}

\usepackage[top=37mm,bottom=37mm,left=27mm,right=27mm]{geometry}
\sloppy
\pagestyle{myheadings}
\usepackage{amssymb}
\usepackage[utf8]{inputenc}
\usepackage{makeidx}
\makeindex
\usepackage{color}
\definecolor{FireBrick}{rgb}{0.5812,0.0074,0.0083}
\definecolor{RoyalBlue}{rgb}{0.0236,0.0894,0.6179}
\definecolor{RoyalGreen}{rgb}{0.0236,0.6179,0.0894}
\definecolor{RoyalRed}{rgb}{0.6179,0.0236,0.0894}
\definecolor{LightBlue}{rgb}{0.8544,0.9511,1.0000}
\definecolor{Black}{rgb}{0.0,0.0,0.0}

\definecolor{linkColor}{rgb}{0.0,0.0,0.554}
\definecolor{citeColor}{rgb}{0.0,0.0,0.554}
\definecolor{fileColor}{rgb}{0.0,0.0,0.554}
\definecolor{urlColor}{rgb}{0.0,0.0,0.554}
\definecolor{promptColor}{rgb}{0.0,0.0,0.589}
\definecolor{brkpromptColor}{rgb}{0.589,0.0,0.0}
\definecolor{gapinputColor}{rgb}{0.589,0.0,0.0}
\definecolor{gapoutputColor}{rgb}{0.0,0.0,0.0}

%%  for a long time these were red and blue by default,
%%  now black, but keep variables to overwrite
\definecolor{FuncColor}{rgb}{0.0,0.0,0.0}
%% strange name because of pdflatex bug:
\definecolor{Chapter }{rgb}{0.0,0.0,0.0}
\definecolor{DarkOlive}{rgb}{0.1047,0.2412,0.0064}


\usepackage{fancyvrb}

\usepackage{mathptmx,helvet}
\usepackage[T1]{fontenc}
\usepackage{textcomp}


\usepackage[
            pdftex=true,
            bookmarks=true,        
            a4paper=true,
            pdftitle={Written with GAPDoc},
            pdfcreator={LaTeX with hyperref package / GAPDoc},
            colorlinks=true,
            backref=page,
            breaklinks=true,
            linkcolor=linkColor,
            citecolor=citeColor,
            filecolor=fileColor,
            urlcolor=urlColor,
            pdfpagemode={UseNone}, 
           ]{hyperref}

\newcommand{\maintitlesize}{\fontsize{50}{55}\selectfont}

% write page numbers to a .pnr log file for online help
\newwrite\pagenrlog
\immediate\openout\pagenrlog =\jobname.pnr
\immediate\write\pagenrlog{PAGENRS := [}
\newcommand{\logpage}[1]{\protect\write\pagenrlog{#1, \thepage,}}
%% were never documented, give conflicts with some additional packages

\newcommand{\GAP}{\textsf{GAP}}

%% nicer description environments, allows long labels
\usepackage{enumitem}
\setdescription{style=nextline}

%% depth of toc
\setcounter{tocdepth}{1}





%% command for ColorPrompt style examples
\newcommand{\gapprompt}[1]{\color{promptColor}{\bfseries #1}}
\newcommand{\gapbrkprompt}[1]{\color{brkpromptColor}{\bfseries #1}}
\newcommand{\gapinput}[1]{\color{gapinputColor}{#1}}


\begin{document}

\logpage{[ 0, 0, 0 ]}
\begin{titlepage}
\mbox{}\vfill

\begin{center}{\maintitlesize \textbf{ lpres \\
\mbox{}}}\\
\vfill

\hypersetup{pdftitle={ lpres }}
\markright{\scriptsize \mbox{}\hfill  lpres  \hfill\mbox{}}
{\Huge \textbf{ Nilpotent Quotients of L\texttt{\symbol{45}}Presented Groups \\
\mbox{}}}\\
\vfill

{\Huge  1.1.2 \mbox{}}\\[1cm]
{ 22 May 2026 \mbox{}}\\[1cm]
\mbox{}\\[2cm]
{\Large \textbf{\strut  Ren{\a'e} Hartung  \strut\mbox{}}}\\
\hypersetup{pdfauthor={ Ren{\a'e} Hartung }}
\end{center}\vfill

\mbox{}\\
\end{titlepage}

\newpage\setcounter{page}{2}
\newpage

\def\contentsname{Contents\logpage{[ 0, 0, 1 ]}}

\tableofcontents
\newpage

 
\chapter{\textcolor{Chapter }{The \textsf{lpres} package}}\logpage{[ 1, 0, 0 ]}
\hyperdef{L}{X86B8787287B59CA4}{}
{
 This package was written by Ren{\a'e} Hartung in 2009; maintenance has been
overtaken by Laurent Bartholdi, who translated the manual to \textsf{GAPDoc}. 
\section{\textcolor{Chapter }{Introduction}}\logpage{[ 1, 1, 0 ]}
\hyperdef{L}{X7DFB63A97E67C0A1}{}
{
 In 1980, Grigorchuk \cite{Grigorchuk80} gave an example of an infinite, finitely generated torsion group which
provided a first explicit counter\texttt{\symbol{45}}example to the General
Burnside Problem. This counter\texttt{\symbol{45}}example is nowadays called
the \emph{Grigorchuk group} and was originally defined as a group of transformations of the unit interval
which preserve the Lebesgue measure. Beside being a
counter\texttt{\symbol{45}}example to the General Burnside Problem, the
Grigorchuk group was a first example of a group with an intermediate growth
function (see \cite{Grigorchuk83}) and was used in the construction of a finitely presented amenable group
which is not elementary amenable (see \cite{Grigorchuk98}).

 The Grigorchuk group is not finitely presentable (see \cite{Grigorchuk99}). However, in 1985, Igor Lysenok (see \cite{Lysenok85}) determined the following recursive presentation for the Grigorchuk group: 
\[\langle a,b,c,d\mid
a^2,b^2,c^2,d^2,bcd,[d,d^a]^{\sigma^n},[d,d^{acaca}]^{\sigma^n},
(n\in{\ensuremath{\mathbb N}})\rangle,\]
 where $\sigma$ is the homomorphism of the free group over $\{a,b,c,d\}$ which is induced by $a\mapsto c^a, b\mapsto d, c\mapsto b$, and $d\mapsto c$. Hence, the infinitely many relators of this recursive presentation can be
described in finite terms using powers of the endomorphism $\sigma$.

 In 2003, Bartholdi \cite{Bartholdi03} introduced the notion of an \emph{L\texttt{\symbol{45}}presentation} for presentations of this type; that is, a group presentation of the form 
\[G=\left\langle S \left| Q\cup \bigcup_{\varphi\in\Phi^*}
R^\varphi\right.\right\rangle,\]
 where $\Phi^*$ denotes the free monoid generated by a set of free group endomorphisms $\Phi$. He proved that various branch groups are finitely
L\texttt{\symbol{45}}presented but not finitely presentable and that every
free group in a variety of groups satisfying finitely many identities is
finitely L\texttt{\symbol{45}}presented (e.g. the Free
Burnside\texttt{\symbol{45}} and the Free $n$\texttt{\symbol{45}}Engel groups).

 The \textsf{lpres}\texttt{\symbol{45}}package defines new \textsf{GAP} objects to work with finitely L\texttt{\symbol{45}}presented groups. The main
part of the package is a nilpotent quotient algorithm for finitely
L\texttt{\symbol{45}}presented groups; that is, an algorithm which takes as
input a finitelyL\texttt{\symbol{45}}presented group $G$ and a positive integer $c$. It computes a polycyclic presentation for the lower central series quotient $G/\gamma_{c+1}(G)$. Therefore, a nilpotent quotient algorithm can be used to determine the
abelian invariants of the lower central series sections $\gamma_c(G)/\gamma_{c+1}(G)$ and the largest nilpotent quotient of $G$ if it exists.

 Our nilpotent quotient algorithm generalizes Nickel's algorithm for finitely
presented groups (see \cite{Nickel96}) which is implemented in the \textsf{NQ}\texttt{\symbol{45}}package; see \cite{nq}. In difference to the \textsf{NQ}\texttt{\symbol{45}}package, the \textsf{lpres}\texttt{\symbol{45}}package is implemented in \textsf{GAP} only.

 Since finite L\texttt{\symbol{45}}presentations generalize finite
presentations, our algorithm also applies to finitely presented groups. It
coincides with Nickel's algorithm in this special case.

 A detailed description of our algorithm can be found in \cite{BEH08} or in the diploma thesis \cite{H08}. Furthermore the \textsf{lpres}\texttt{\symbol{45}}package includes the algorithms of \cite{MR2678952} for approximating the Schur multiplier of finitely
L\texttt{\symbol{45}}presented groups. 

 The \textsf{lpres}\texttt{\symbol{45}}package also includes the
Reidemeister\texttt{\symbol{45}}Schreier algorithm from \cite{MR2876891}. L\texttt{\symbol{45}}presented groups were introduced as a tool to
understand self\texttt{\symbol{45}}similar groups such as the Grigorchuk
group. As such, \textsf{lpres} works in close contact with the package \textsf{fr}. See \cite{MR3133711} for more on the relationships between L\texttt{\symbol{45}}presented groups
and self\texttt{\symbol{45}}similar groups.

 Finally, we note that we use the term "algorithm" somewhat loosely: many of
the algorithms in this package are in fact semi\texttt{\symbol{45}}algorithms,
guaranteed to give a correct answer but not guaranteed to terminate. }

 }

 
\chapter{\textcolor{Chapter }{An Introduction to L\texttt{\symbol{45}}presented groups}}\logpage{[ 2, 0, 0 ]}
\hyperdef{L}{X7AEB47327D75B633}{}
{
 
\section{\textcolor{Chapter }{Definitions}}\logpage{[ 2, 1, 0 ]}
\hyperdef{L}{X84541F61810C741D}{}
{
 Let $S$ be an alphabet, $Q$ and $R$ be subsets of the free group $F_S$ over this alphabet, and $\Phi$ be a set of free group endomorphisms $\varphi\colon F_S\to F_S$. An \emph{L\texttt{\symbol{45}}presentation} is a quadruple $(S,Q,\Phi,R)$ and it is called \emph{finite} if the sets $S$, $Q$, $\Phi$, and $R$ are finite. A (finite) L\texttt{\symbol{45}}presentation $(S,Q,\Phi,R)$ defines the (\emph{finitely}) \emph{L\texttt{\symbol{45}}presented group} 
\[ G=\left\langle S \left| Q\cup
\bigcup_{\varphi\in\Phi^*}R^\varphi\right.\right\rangle\]
 where $\Phi^*$ denotes the free monoid generated by $\Phi$; that is, the closure of $\Phi\cup\{\rm id\}$ under composition.

 The elements in $Q$ are the \emph{fixed relators} and the elements in $R$ are the \emph{iterated relators} of the L\texttt{\symbol{45}}presentation $(S,Q,\Phi,R)$. An L\texttt{\symbol{45}}presentation of the form $(S,\emptyset,\Phi,R)$ is an \emph{ascending L\texttt{\symbol{45}}presentation} and it is an \emph{invariant L\texttt{\symbol{45}}presentation} if the normal subgroup 
\[K=\left\langle Q\cup \bigcup_{\varphi\in\Phi^*}R^\varphi\right\rangle^{F_S}\]
 is $\varphi$\texttt{\symbol{45}}invariant for each $\varphi\in\Phi$; that is, if $K$ satisfies $K^\varphi\subset K$ for each $\varphi\in\Phi$. Note that every ascending L\texttt{\symbol{45}}presentation is invariant and
for each L\texttt{\symbol{45}}presentation $(S,Q,\Phi,R)$ there is a unique \emph{underlying ascending L\texttt{\symbol{45}}presentation} $(S,\emptyset,\Phi,R)$ which is invariant. In general it is not decidable whether or not a given
L\texttt{\symbol{45}}presentation is invariant as this would require a
solution to the word\texttt{\symbol{45}}problem.

 In the remainder of this manual, an L\texttt{\symbol{45}}presented group is
always finitely L\texttt{\symbol{45}}presented.

 }

 
\section{\textcolor{Chapter }{Creating an L\texttt{\symbol{45}}presented group}}\label{createlp}
\logpage{[ 2, 2, 0 ]}
\hyperdef{L}{X81065E797A486D0F}{}
{
 The construction of an L\texttt{\symbol{45}}presented group is similar to the
construction of a finitely presented group (see Chapter  (\textbf{Reference: Finitely Presented Groups}) of the \textsf{GAP} Reference manual for further details). 

\subsection{\textcolor{Chapter }{LPresentedGroup}}
\logpage{[ 2, 2, 1 ]}\nobreak
\hyperdef{L}{X7BBBE4C082AE4D5A}{}
{\noindent\textcolor{FuncColor}{$\triangleright$\enspace\texttt{LPresentedGroup({\mdseries\slshape F, frels, endos, irels})\index{LPresentedGroup@\texttt{LPresentedGroup}}
\label{LPresentedGroup}
}\hfill{\scriptsize (function)}}\\


 returns the \textsf{GAP} object of an L\texttt{\symbol{45}}presented group with the underlying free
group \mbox{\texttt{\mdseries\slshape F}}, the fixed relators \mbox{\texttt{\mdseries\slshape frels}}, the set of endomorphisms \mbox{\texttt{\mdseries\slshape endos}}, and the iterated relators \mbox{\texttt{\mdseries\slshape irels}}. The input variables \mbox{\texttt{\mdseries\slshape frels}} and \mbox{\texttt{\mdseries\slshape irels}} need to be finite subsets of the underlying free group \mbox{\texttt{\mdseries\slshape F}} and \mbox{\texttt{\mdseries\slshape endos}} needs to be a finite list of homomorphisms $F\to F$.

 For example, the Grigorchuk group, 
\[ \Big\langle a,b,c,d \Big|
a^2,b^2,c^2,d^2,bcd,[d,d^a]^{\sigma^n},[d,d^{acaca}]^{\sigma^n},(n\in{\ensuremath{\mathbb
N}}_0) \Big\rangle,\]
 can be constructed as follows. 
\begin{Verbatim}[commandchars=!@|,fontsize=\small,frame=single,label=Example]
  !gapprompt@gap>| !gapinput@F:=FreeGroup( "a", "b", "c", "d" );|
  <free group on the generators [ a, b, c, d ]>
  !gapprompt@gap>| !gapinput@AssignGeneratorVariables( F );|
  #I  Assigned the global variables [ a, b, c, d ]
  !gapprompt@gap>| !gapinput@frels:=[a^2, b^2, c^2, d^2, b*c*d];;|
  !gapprompt@gap>| !gapinput@endos:=[GroupHomomorphismByImagesNC( F, F, [a, b, c, d], [c^a, d, b, c])];;|
  !gapprompt@gap>| !gapinput@irels:=[Comm( d, d^a ), Comm( d, d^(a*c*a*c*a) )];;|
  !gapprompt@gap>| !gapinput@G:=LPresentedGroup( F, frels, endos, irels );|
  <L-presented group on the generators [ a, b, c, d ]>
\end{Verbatim}
 }

 There are various examples of finitely L\texttt{\symbol{45}}presented groups
available in the library of the \textsf{lpres}\texttt{\symbol{45}}package. 

\subsection{\textcolor{Chapter }{ExamplesOfLPresentations}}
\logpage{[ 2, 2, 2 ]}\nobreak
\hyperdef{L}{X79A034B8851444C9}{}
{\noindent\textcolor{FuncColor}{$\triangleright$\enspace\texttt{ExamplesOfLPresentations({\mdseries\slshape n})\index{ExamplesOfLPresentations@\texttt{ExamplesOfLPresentations}}
\label{ExamplesOfLPresentations}
}\hfill{\scriptsize (function)}}\\


 returns some well\texttt{\symbol{45}}known examples of finitely
L\texttt{\symbol{45}}presented groups. The input of this function needs to be
a positive integer at most $10$.

 
\begin{description}
\item[{n=1}] The Grigorchuk group on 4 generators; cf. \cite{Grigorchuk80}, \cite{Lysenok85}, and \cite[Theorem 4.6]{Bartholdi03},
\item[{n=2}] the Grigorchuk group on 3 generators; cf. \cite{Grigorchuk80}, \cite{Lysenok85}, and \cite[Theorem 4.6]{Bartholdi03},
\item[{n=3}] the lamplighter group ${\ensuremath{\mathbb Z}}/2\wr{\ensuremath{\mathbb Z}}$; cf. \cite[Theorem 4.1]{Bartholdi03},
\item[{n=4}] the Brunner\texttt{\symbol{45}}Sidki\texttt{\symbol{45}}Vieira group; cf. \cite{BrunnerVieiraSidki99} and \cite[Theorem 4.4]{Bartholdi03},
\item[{n=5}] the Grigorchuk supergroup; cf. \cite{BartholdiGrigorchuk02} and \cite[Theorem 4.6]{Bartholdi03},
\item[{n=6}] the Fabrykowski\texttt{\symbol{45}}Gupta group; cf. \cite{FabrykowskiGupta85} and \cite{BEH08},
\item[{n=7}] the Gupta\texttt{\symbol{45}}Sidki group; cf. \cite{Sidki87} and \cite{BEH08},
\item[{n=8}] an index\texttt{\symbol{45}}$3$ subgroup of the Gupta\texttt{\symbol{45}}Sidki group,
\item[{n=9}] the Basilica group; cf. \cite{GrigorchukZuk02} and \cite{BartholdiVirag05},
\item[{n=10}] Baumslag's finitely generated, infinitely related group with a trivial
multiplier; cf. \cite{Baumslag71}.
\end{description}
 }

 Furthermore, every free group in a variety of groups satisfying finitely many
identities is finitely L\texttt{\symbol{45}}presented. Some of these groups
are available from the \textsf{lpres}\texttt{\symbol{45}}package using the following operations; for further
details we refer to the diploma thesis \cite{H08}. 

\subsection{\textcolor{Chapter }{FreeEngelGroup}}
\logpage{[ 2, 2, 3 ]}\nobreak
\hyperdef{L}{X7DA323A87E7B6A7C}{}
{\noindent\textcolor{FuncColor}{$\triangleright$\enspace\texttt{FreeEngelGroup({\mdseries\slshape num, n})\index{FreeEngelGroup@\texttt{FreeEngelGroup}}
\label{FreeEngelGroup}
}\hfill{\scriptsize (operation)}}\\


 returns an L\texttt{\symbol{45}}presentation for the free \mbox{\texttt{\mdseries\slshape n}}\texttt{\symbol{45}}Engel group on \mbox{\texttt{\mdseries\slshape num}} generators; that is, the free group in the variety of \mbox{\texttt{\mdseries\slshape num}}\texttt{\symbol{45}}generated groups satisfying the \mbox{\texttt{\mdseries\slshape n}}\texttt{\symbol{45}}Engel identity. }

 

\subsection{\textcolor{Chapter }{FreeBurnsideGroup}}
\logpage{[ 2, 2, 4 ]}\nobreak
\hyperdef{L}{X81C3537083E40A5C}{}
{\noindent\textcolor{FuncColor}{$\triangleright$\enspace\texttt{FreeBurnsideGroup({\mdseries\slshape num, exp})\index{FreeBurnsideGroup@\texttt{FreeBurnsideGroup}}
\label{FreeBurnsideGroup}
}\hfill{\scriptsize (operation)}}\\


 returns an L\texttt{\symbol{45}}presentation for the free Burnside group on \mbox{\texttt{\mdseries\slshape num}} generators with exponent \mbox{\texttt{\mdseries\slshape exp}}; that is, the free group on \mbox{\texttt{\mdseries\slshape num}} generators in the variety of groups with exponent \mbox{\texttt{\mdseries\slshape exp}}. }

 

\subsection{\textcolor{Chapter }{FreeNilpotentGroup}}
\logpage{[ 2, 2, 5 ]}\nobreak
\hyperdef{L}{X8796306C7A7924D1}{}
{\noindent\textcolor{FuncColor}{$\triangleright$\enspace\texttt{FreeNilpotentGroup({\mdseries\slshape num, c})\index{FreeNilpotentGroup@\texttt{FreeNilpotentGroup}}
\label{FreeNilpotentGroup}
}\hfill{\scriptsize (operation)}}\\


 returns an L\texttt{\symbol{45}}presentation for the free nilpotent group of
class \mbox{\texttt{\mdseries\slshape c}} on \mbox{\texttt{\mdseries\slshape num}} generators; that is, the free group in the variety of \mbox{\texttt{\mdseries\slshape num}}\texttt{\symbol{45}}generated, nilpotent groups with nilpotency class \mbox{\texttt{\mdseries\slshape c}}. }

 

\subsection{\textcolor{Chapter }{GeneralizedFabrykowskiGuptaLpGroup}}
\logpage{[ 2, 2, 6 ]}\nobreak
\hyperdef{L}{X81450ABA81F0FCE5}{}
{\noindent\textcolor{FuncColor}{$\triangleright$\enspace\texttt{GeneralizedFabrykowskiGuptaLpGroup({\mdseries\slshape n})\index{GeneralizedFabrykowskiGuptaLpGroup@\texttt{GeneralizedFabrykowskiGuptaLpGroup}}
\label{GeneralizedFabrykowskiGuptaLpGroup}
}\hfill{\scriptsize (operation)}}\\


 returns an L\texttt{\symbol{45}}presentation for the \mbox{\texttt{\mdseries\slshape n}}\texttt{\symbol{45}}th generalized Fabrykowski\texttt{\symbol{45}}Gupta group
as constructed in \cite{BEH08}. }

 

\subsection{\textcolor{Chapter }{LamplighterGroup (llint)}}
\logpage{[ 2, 2, 7 ]}\nobreak
\hyperdef{L}{X83BF8C597E1DC266}{}
{\noindent\textcolor{FuncColor}{$\triangleright$\enspace\texttt{LamplighterGroup({\mdseries\slshape filter, int})\index{LamplighterGroup@\texttt{LamplighterGroup}!llint}
\label{LamplighterGroup:llint}
}\hfill{\scriptsize (operation)}}\\
\noindent\textcolor{FuncColor}{$\triangleright$\enspace\texttt{LamplighterGroup({\mdseries\slshape filter, pcgroup})\index{LamplighterGroup@\texttt{LamplighterGroup}!llpcgroup}
\label{LamplighterGroup:llpcgroup}
}\hfill{\scriptsize (operation)}}\\


 returns a finite L\texttt{\symbol{45}}presentation for the lamplighter group
on \mbox{\texttt{\mdseries\slshape int}} lamp states in the first case, if \mbox{\texttt{\mdseries\slshape filter}} is the filter \texttt{IsLpGroup}. In the second case, the group \mbox{\texttt{\mdseries\slshape pcgroup}} must be a finite cyclic group. Then the method returns a finite
L\texttt{\symbol{45}}presentation for the lamplighter group on \texttt{Size(pcgroup)} lamp states; for details on the L\texttt{\symbol{45}}presentation see \cite{Bartholdi03}. 
\begin{Verbatim}[commandchars=!@|,fontsize=\small,frame=single,label=Example]
  !gapprompt@gap>| !gapinput@LamplighterGroup( IsLpGroup, 2 );|
  <L-presented group on the generators [ a, t, u ]>
  !gapprompt@gap>| !gapinput@LamplighterGroup( IsLpGroup, CyclicGroup(3) );|
  <L-presented group on the generators [ a, t, u ]>
\end{Verbatim}
}

 

\subsection{\textcolor{Chapter }{EmbeddingOfIASubgroup}}
\logpage{[ 2, 2, 8 ]}\nobreak
\hyperdef{L}{X7DBA63A37853BE46}{}
{\noindent\textcolor{FuncColor}{$\triangleright$\enspace\texttt{EmbeddingOfIASubgroup({\mdseries\slshape a})\index{EmbeddingOfIASubgroup@\texttt{EmbeddingOfIASubgroup}}
\label{EmbeddingOfIASubgroup}
}\hfill{\scriptsize (operation)}}\\


 computes an L\texttt{\symbol{45}}presentation for the
IA\texttt{\symbol{45}}automorphism group of a free group. This is the subgroup
of automorphisms of a free group $f$ that act trivially on the abelianization of $f$.

 The L\texttt{\symbol{45}}presentation is taken from \cite{DayPutman}. 
\begin{Verbatim}[commandchars=!@|,fontsize=\small,frame=single,label=Example]
  !gapprompt@gap>| !gapinput@f := FreeGroup(3);|
  <free group on the generators [ f1, f2, f3 ]>
  !gapprompt@gap>| !gapinput@a := AutomorphismGroup(f);|
  <group of size infinity with 3 generators>
  !gapprompt@gap>| !gapinput@ia := Source(EmbeddingOfIASubgroup(a));|
  <invariant LpGroup on the generators [ C(1,2), C(1,3), C(2,1), C(2,3), C(3,1), C(3,2), M(1,[2,3]),
    M(2,[1,3]), M(3,[1,2]) ]>
  !gapprompt@gap>| !gapinput@rank := 3;|
  3
  !gapprompt@gap>| !gapinput@q := NilpotentQuotient(ia,rank);;|
  !gapprompt@gap>| !gapinput@lcs := LowerCentralSeries(q);;|
  !gapprompt@gap>| !gapinput@for i in [1..Length(lcs)-1] do|
  !gapprompt@>| !gapinput@        r := AbelianInvariants(lcs[i]/lcs[i+1]);|
  !gapprompt@>| !gapinput@        Print(i); if i>3 then Print("th"); else Print(ELM_LIST(["st","nd","rd"],i)); fi;|
  !gapprompt@>| !gapinput@        Print(" quotient: abelian invariants ",r," (collected ",Collected(r),")\n");|
  !gapprompt@>| !gapinput@    od;|
  1st quotient: abelian invariants [ 0, 0, 0, 0, 0, 0, 0, 0, 0 ] (collected [ [ 0, 9 ] ])
  2nd quotient: abelian invariants [ 0, 0, 0, 0, 0, 0, 0, 0, 0, 0, 0, 0, 0, 0, 0, 0, 0, 0
   ] (collected [ [ 0, 18 ] ])
  3rd quotient: abelian invariants [ 0, 0, 0, 0, 0, 0, 0, 0, 0, 0, 0, 0, 0, 0, 0, 0, 0, 0, 0, 0, 0,
    0, 0, 0, 0, 0, 0, 0, 0, 0, 0, 0, 0, 0, 0, 0, 0, 0, 0, 0, 0, 0, 0, 2, 2, 2, 2, 2, 2, 2, 2, 2, 2,
    2, 2, 2, 2, 3, 3, 3, 3, 3, 3, 3, 3, 3 ] (collected [ [ 0, 43 ], [ 2, 14 ], [ 3, 9 ] ])
\end{Verbatim}
 }

 }

 
\section{\textcolor{Chapter }{The underlying free group}}\logpage{[ 2, 3, 0 ]}
\hyperdef{L}{X80B65AF48662DE70}{}
{
 An L\texttt{\symbol{45}}presented group is defined as an image of its
underlying free group. Note that these are two different \textsf{GAP} objects. The elements of the L\texttt{\symbol{45}}presented group are
represented by words in the underlying free group. 

\subsection{\textcolor{Chapter }{FreeGroupOfLpGroup}}
\logpage{[ 2, 3, 1 ]}\nobreak
\hyperdef{L}{X7F883CC57A3CCAC7}{}
{\noindent\textcolor{FuncColor}{$\triangleright$\enspace\texttt{FreeGroupOfLpGroup({\mdseries\slshape lpgroup})\index{FreeGroupOfLpGroup@\texttt{FreeGroupOfLpGroup}}
\label{FreeGroupOfLpGroup}
}\hfill{\scriptsize (attribute)}}\\
\textbf{\indent Returns:\ }
the underlying free group of the L\texttt{\symbol{45}}presented group \mbox{\texttt{\mdseries\slshape lpgroup}}

}

 

\subsection{\textcolor{Chapter }{FreeGeneratorsOfLpGroup}}
\logpage{[ 2, 3, 2 ]}\nobreak
\hyperdef{L}{X838079A587E8CF43}{}
{\noindent\textcolor{FuncColor}{$\triangleright$\enspace\texttt{FreeGeneratorsOfLpGroup({\mdseries\slshape lpgroup})\index{FreeGeneratorsOfLpGroup@\texttt{FreeGeneratorsOfLpGroup}}
\label{FreeGeneratorsOfLpGroup}
}\hfill{\scriptsize (attribute)}}\\
\textbf{\indent Returns:\ }
the generators of the free group which underlies the
L\texttt{\symbol{45}}presented group \mbox{\texttt{\mdseries\slshape lpgroup}}

}

 

\subsection{\textcolor{Chapter }{GeneratorsOfGroup}}
\logpage{[ 2, 3, 3 ]}\nobreak
\hyperdef{L}{X79C44528864044C5}{}
{\noindent\textcolor{FuncColor}{$\triangleright$\enspace\texttt{GeneratorsOfGroup({\mdseries\slshape lpgroup})\index{GeneratorsOfGroup@\texttt{GeneratorsOfGroup}}
\label{GeneratorsOfGroup}
}\hfill{\scriptsize (attribute)}}\\
\textbf{\indent Returns:\ }
the generators of the L\texttt{\symbol{45}}presented group \mbox{\texttt{\mdseries\slshape lpgroup}}. These are the images of the generators of the underlying free group under
the natural homomorphism.

}

 

\subsection{\textcolor{Chapter }{UnderlyingElement}}
\logpage{[ 2, 3, 4 ]}\nobreak
\hyperdef{L}{X85C405D57F65048A}{}
{\noindent\textcolor{FuncColor}{$\triangleright$\enspace\texttt{UnderlyingElement({\mdseries\slshape elm})\index{UnderlyingElement@\texttt{UnderlyingElement}}
\label{UnderlyingElement}
}\hfill{\scriptsize (operation)}}\\


 returns the preimage of an L\texttt{\symbol{45}}presented group element \mbox{\texttt{\mdseries\slshape elm}} in the underlying free group. More precisely, each element of an
L\texttt{\symbol{45}}presented group is represented by an element in the free
group. This method returns the corresponding element in the free group. }

 

\subsection{\textcolor{Chapter }{ElementOfLpGroup}}
\logpage{[ 2, 3, 5 ]}\nobreak
\hyperdef{L}{X8573CDF57CB216D7}{}
{\noindent\textcolor{FuncColor}{$\triangleright$\enspace\texttt{ElementOfLpGroup({\mdseries\slshape fam, elm})\index{ElementOfLpGroup@\texttt{ElementOfLpGroup}}
\label{ElementOfLpGroup}
}\hfill{\scriptsize (operation)}}\\


 returns the element in the L\texttt{\symbol{45}}presented group represented by
the word \mbox{\texttt{\mdseries\slshape elm}} on the generators of the underlying free group, if \mbox{\texttt{\mdseries\slshape fam}} is the family of L\texttt{\symbol{45}}presented group elements. 
\begin{Verbatim}[commandchars=!@|,fontsize=\small,frame=single,label=Example]
  !gapprompt@gap>| !gapinput@F:=FreeGroup( 2 );;|
  !gapprompt@gap>| !gapinput@G:=LPresentedGroup( F, [ F.1^2 ], [ IdentityMapping( F ) ], [ F.2 ] );;|
  !gapprompt@gap>| !gapinput@FreeGroupOfLpGroup( G ) = F;|
  true
  !gapprompt@gap>| !gapinput@GeneratorsOfGroup( G );|
  [ f1, f2 ]
  !gapprompt@gap>| !gapinput@FreeGeneratorsOfLpGroup( G );|
  [ f1, f2 ]
  !gapprompt@gap>| !gapinput@last = last2;|
  false
  !gapprompt@gap>| !gapinput@UnderlyingElement( G.1 );|
  f1
  !gapprompt@gap>| !gapinput@last in F;|
  true
  !gapprompt@gap>| !gapinput@ElementOfLpGroup( ElementsFamily( FamilyObj( G ) ), last2 ) in G;|
  true
\end{Verbatim}
 }

 }

 
\section{\textcolor{Chapter }{Accessing an L\texttt{\symbol{45}}presentation}}\logpage{[ 2, 4, 0 ]}
\hyperdef{L}{X847047F083826C00}{}
{
 The fixed relators, the iterated relators, and the endomorphisms of an
L\texttt{\symbol{45}}presented group are accessible with the following
methods. 

\subsection{\textcolor{Chapter }{FixedRelatorsOfLpGroup}}
\logpage{[ 2, 4, 1 ]}\nobreak
\hyperdef{L}{X7CD9BE57815552FF}{}
{\noindent\textcolor{FuncColor}{$\triangleright$\enspace\texttt{FixedRelatorsOfLpGroup({\mdseries\slshape lpgroup})\index{FixedRelatorsOfLpGroup@\texttt{FixedRelatorsOfLpGroup}}
\label{FixedRelatorsOfLpGroup}
}\hfill{\scriptsize (attribute)}}\\
\textbf{\indent Returns:\ }
the fixed relators of the L\texttt{\symbol{45}}presented group \mbox{\texttt{\mdseries\slshape lpgroup}} as elements of the underlying free group.

}

 

\subsection{\textcolor{Chapter }{IteratedRelatorsOfLpGroup}}
\logpage{[ 2, 4, 2 ]}\nobreak
\hyperdef{L}{X7C468D1C81964268}{}
{\noindent\textcolor{FuncColor}{$\triangleright$\enspace\texttt{IteratedRelatorsOfLpGroup({\mdseries\slshape lpgroup})\index{IteratedRelatorsOfLpGroup@\texttt{IteratedRelatorsOfLpGroup}}
\label{IteratedRelatorsOfLpGroup}
}\hfill{\scriptsize (attribute)}}\\
\textbf{\indent Returns:\ }
the iterated relators of the L\texttt{\symbol{45}}presented group \mbox{\texttt{\mdseries\slshape lpgroup}} as elements of the underlying free group.

}

 

\subsection{\textcolor{Chapter }{EndomorphismsOfLpGroup}}
\logpage{[ 2, 4, 3 ]}\nobreak
\hyperdef{L}{X85D253888263A3F6}{}
{\noindent\textcolor{FuncColor}{$\triangleright$\enspace\texttt{EndomorphismsOfLpGroup({\mdseries\slshape lpgroup})\index{EndomorphismsOfLpGroup@\texttt{EndomorphismsOfLpGroup}}
\label{EndomorphismsOfLpGroup}
}\hfill{\scriptsize (attribute)}}\\
\textbf{\indent Returns:\ }
the endomorphisms of the L\texttt{\symbol{45}}presented group \mbox{\texttt{\mdseries\slshape lpgroup}} as endomorphisms of the underlying free group.



 
\begin{Verbatim}[commandchars=!@|,fontsize=\small,frame=single,label=Example]
  !gapprompt@gap>| !gapinput@F:=FreeGroup( 2 );;|
  !gapprompt@gap>| !gapinput@G:=LPresentedGroup( F, [ F.1^2 ], [ IdentityMapping( F ) ], [ F.2 ] );|
  <L-presented group on the generators [ f1, f2 ]>
  !gapprompt@gap>| !gapinput@FixedRelatorsOfLpGroup( G );|
  [ f1^2 ]
  !gapprompt@gap>| !gapinput@IteratedRelatorsOfLpGroup( G );|
  [ f2 ]
  !gapprompt@gap>| !gapinput@EndomorphismsOfLpGroup( G );|
  [ IdentityMapping( <free group on the generators [ f1, f2 ]> ) ]
\end{Verbatim}
 }

 }

 
\section{\textcolor{Chapter }{Attributes and properties of L\texttt{\symbol{45}}presented groups}}\logpage{[ 2, 5, 0 ]}
\hyperdef{L}{X817DA8E686311B54}{}
{
 For the method\texttt{\symbol{45}}selection of the nilpotent quotient
algorithm, an L\texttt{\symbol{45}}presented group may have the following
attributes and properties. 

\subsection{\textcolor{Chapter }{UnderlyingAscendingLPresentation}}
\logpage{[ 2, 5, 1 ]}\nobreak
\hyperdef{L}{X85E77B29796AB730}{}
{\noindent\textcolor{FuncColor}{$\triangleright$\enspace\texttt{UnderlyingAscendingLPresentation({\mdseries\slshape lpgroup})\index{UnderlyingAscendingLPresentation@\texttt{UnderlyingAscendingLPresentation}}
\label{UnderlyingAscendingLPresentation}
}\hfill{\scriptsize (attribute)}}\\


 returns the underlying ascending L\texttt{\symbol{45}}presentation of \mbox{\texttt{\mdseries\slshape lpgroup}}; that is, if \mbox{\texttt{\mdseries\slshape lpgroup}} is finitely L\texttt{\symbol{45}}presented by $(S,Q,\Phi,R)$, the underlying ascending L\texttt{\symbol{45}}presentation is $(S,\emptyset,\Phi,R)$.

 }

 

\subsection{\textcolor{Chapter }{UnderlyingInvariantLPresentation}}
\logpage{[ 2, 5, 2 ]}\nobreak
\hyperdef{L}{X86F017E085082624}{}
{\noindent\textcolor{FuncColor}{$\triangleright$\enspace\texttt{UnderlyingInvariantLPresentation({\mdseries\slshape lpgroup})\index{UnderlyingInvariantLPresentation@\texttt{UnderlyingInvariantLPresentation}}
\label{UnderlyingInvariantLPresentation}
}\hfill{\scriptsize (attribute)}}\\


 attempts to compute a ``good'' underlying invariant
L\texttt{\symbol{45}}presentation for \mbox{\texttt{\mdseries\slshape lpgroup}}; that is, if \mbox{\texttt{\mdseries\slshape lpgroup}} is finitely L\texttt{\symbol{45}}presented by $(S,Q,\Phi,R)$, then this method seeks to find a subset $Q'\subseteq Q$ such that $(S,Q',\Phi,R)$ is an invariant L\texttt{\symbol{45}}presentation. Note that there is always
the underlying ascending L\texttt{\symbol{45}}presentation $(S,\emptyset,\Phi,R)$. However, for the efficiency of the nilpotent quotient algorithm it is
important that the subset $Q'$ is as big as possible.

 Since it is undecidable, in general, whether or not a given
L\texttt{\symbol{45}}presentation is invariant, there is no algorithm which
can determine the best possible underlying invariant
L\texttt{\symbol{45}}presentation. The method implemented for this attribute
tries to compute a ``good'' invariant L\texttt{\symbol{45}}presentation and
will return the underlying ascending L\texttt{\symbol{45}}presentation in the
worst case.

 This attribute can be set manually using \texttt{SetUnderlyingInvariantLPresentation}. For instance, the Grigorchuk group 
\[ \Big\langle a,b,c,d \Big| a^2,b^2,c^2,d^2,bcd,[d,d^a]^{\sigma^n},
[d,d^{acaca}]^{\sigma^n},(n\in{\ensuremath{\mathbb N}}_0) \Big\rangle,\]
 is invariantly L\texttt{\symbol{45}}presented and therefore, it should be
constructed as follows: 
\begin{Verbatim}[commandchars=!@|,fontsize=\small,frame=single,label=Example]
  !gapprompt@gap>| !gapinput@F:=FreeGroup( "a", "b", "c", "d" );;|
  !gapprompt@gap>| !gapinput@AssignGeneratorVariables( F );|
  #I  Assigned the global variables [ a, b, c, d ]
  !gapprompt@gap>| !gapinput@frels:=[ a^2, b^2, c^2, d^2, b*c*d ];;|
  !gapprompt@gap>| !gapinput@endos:=[ GroupHomomorphismByImagesNC( F, F, [ a, b, c, d ], [ c^a, d, b, c ]) ];;|
  !gapprompt@gap>| !gapinput@irels:=[ Comm( d, d^a ), Comm( d, d^(a*c*a*c*a) ) ];;|
  !gapprompt@gap>| !gapinput@G:=LPresentedGroup( F, frels, endos, irels );|
  <L-presented group on the generators [ a, b, c, d ]>
  !gapprompt@gap>| !gapinput@SetUnderlyingInvariantLPresentation( G, G );;|
\end{Verbatim}
 }

 

\subsection{\textcolor{Chapter }{IsAscendingLPresentation}}
\logpage{[ 2, 5, 3 ]}\nobreak
\hyperdef{L}{X84E7A9E07A5DFDCF}{}
{\noindent\textcolor{FuncColor}{$\triangleright$\enspace\texttt{IsAscendingLPresentation({\mdseries\slshape lpgroup})\index{IsAscendingLPresentation@\texttt{IsAscendingLPresentation}}
\label{IsAscendingLPresentation}
}\hfill{\scriptsize (property)}}\\


 checks whether the L\texttt{\symbol{45}}presentation of \mbox{\texttt{\mdseries\slshape lpgroup}} is ascending; that is, if the set of fixed relators is empty. This property is
set automatically when creating an L\texttt{\symbol{45}}presented group with
no fixed relators using the function \texttt{LPresentedGroup} (\ref{LPresentedGroup}). }

 

\subsection{\textcolor{Chapter }{IsInvariantLPresentation}}
\logpage{[ 2, 5, 4 ]}\nobreak
\hyperdef{L}{X87F0C52978D99BB5}{}
{\noindent\textcolor{FuncColor}{$\triangleright$\enspace\texttt{IsInvariantLPresentation({\mdseries\slshape lpgroup})\index{IsInvariantLPresentation@\texttt{IsInvariantLPresentation}}
\label{IsInvariantLPresentation}
}\hfill{\scriptsize (property)}}\\


 attempts to check whether the L\texttt{\symbol{45}}presentation of \mbox{\texttt{\mdseries\slshape lpgroup}} is invariant. In general, one cannot decide whether or not a given
L\texttt{\symbol{45}}presentation is invariant. There are mainly two methods
implemented for this property. The first method seeks to find a ``good''
underlying invariant L\texttt{\symbol{45}}presentation using the operation \texttt{UnderlyingInvariantLPresentation} (\ref{UnderlyingInvariantLPresentation}). If this latter L\texttt{\symbol{45}}presentation coincides with the
L\texttt{\symbol{45}}presentation of \mbox{\texttt{\mdseries\slshape lpgroup}}, then \mbox{\texttt{\mdseries\slshape lpgroup}} is invariantly L\texttt{\symbol{45}}presented. If this method fails, then the
second method uses the nilpotent quotient algorithm for
L\texttt{\symbol{45}}presented groups which yields a necessary condition for
an L\texttt{\symbol{45}}presented group to be invariantly
L\texttt{\symbol{45}}presented. Note that the latter method may not terminate.
For instance, both methods fail on Baumslag's finitely generated, infinitely
related group with trivial multiplier returned by \texttt{ExamplesOfLPresentations} (\ref{ExamplesOfLPresentations}). }

 

\subsection{\textcolor{Chapter }{EmbeddingOfAscendingSubgroup}}
\logpage{[ 2, 5, 5 ]}\nobreak
\hyperdef{L}{X783B99E381C5C8BF}{}
{\noindent\textcolor{FuncColor}{$\triangleright$\enspace\texttt{EmbeddingOfAscendingSubgroup({\mdseries\slshape lpgroup})\index{EmbeddingOfAscendingSubgroup@\texttt{EmbeddingOfAscendingSubgroup}}
\label{EmbeddingOfAscendingSubgroup}
}\hfill{\scriptsize (attribute)}}\\


 stores an embedding of an ascendingly L\texttt{\symbol{45}}presented subgroup
of the L\texttt{\symbol{45}}presented group \mbox{\texttt{\mdseries\slshape lpgroup}}. This attribute is set for ascending L\texttt{\symbol{45}}presentations only.
In this case, the identity map of \mbox{\texttt{\mdseries\slshape lpgroup}} is returned. This attribute is used in the \textsf{FR}\texttt{\symbol{45}}package which can construct various finitely
L\texttt{\symbol{45}}presented groups. The embedding is useful for a nilpotent
quotient algorithm of a non\texttt{\symbol{45}}invariantly
L\texttt{\symbol{45}}presented group. }

 }

 
\section{\textcolor{Chapter }{Methods for L\texttt{\symbol{45}}presented groups}}\logpage{[ 2, 6, 0 ]}
\hyperdef{L}{X7B5C48EA7CD8A57E}{}
{
 Some operations are natural extensions of the operations for finitely
generated groups. For example, \texttt{MappedWord(x,gens,imgs)}, when applied to a word \texttt{x} in an L\texttt{\symbol{45}}presented group, returns the group element obtained
by replacing each occurrence of a generator in \texttt{gens} by the corresponding element in the list \texttt{imgs}. The lists \texttt{gens} and \texttt{imgs} need to have the same length.

 Equality test of elements of L\texttt{\symbol{45}}presented groups is
implemented using the operation \texttt{NqEpimorphismNilpotentQuotient} (\textbf{nq: NqEpimorphismNilpotentQuotient}) to compare the images in a nilpotent quotient of the group. The implemented
method successively increases the class of the considered quotient until the
images differ. Hence, this method may not terminate and it will only determine
whether the elements are different.

 

\subsection{\textcolor{Chapter }{EpimorphismFromFpGroup}}
\logpage{[ 2, 6, 1 ]}\nobreak
\hyperdef{L}{X7C81CB1C7F0D7A90}{}
{\noindent\textcolor{FuncColor}{$\triangleright$\enspace\texttt{EpimorphismFromFpGroup({\mdseries\slshape lpgroup, n})\index{EpimorphismFromFpGroup@\texttt{EpimorphismFromFpGroup}}
\label{EpimorphismFromFpGroup}
}\hfill{\scriptsize (operation)}}\\


 returns an epimorphism from a finitely presented group onto \mbox{\texttt{\mdseries\slshape lpgroup}}. The finitely presented group is obtained from \mbox{\texttt{\mdseries\slshape lpgroup}} by applying only words of length at most \mbox{\texttt{\mdseries\slshape n}} in the endomorphisms of \mbox{\texttt{\mdseries\slshape lpgroup}} to the iterated relators of \mbox{\texttt{\mdseries\slshape lpgroup}}. }

 

\subsection{\textcolor{Chapter }{SplitExtensionByAutomorphismsLpGroup}}
\logpage{[ 2, 6, 2 ]}\nobreak
\hyperdef{L}{X7972B0D87EF36536}{}
{\noindent\textcolor{FuncColor}{$\triangleright$\enspace\texttt{SplitExtensionByAutomorphismsLpGroup({\mdseries\slshape lpgroup, H, auts})\index{SplitExtensionByAutomorphismsLpGroup@\texttt{Split}\-\texttt{Extension}\-\texttt{By}\-\texttt{Automorphisms}\-\texttt{Lp}\-\texttt{Group}}
\label{SplitExtensionByAutomorphismsLpGroup}
}\hfill{\scriptsize (operation)}}\\


 returns an L\texttt{\symbol{45}}presentation for the split extension of \mbox{\texttt{\mdseries\slshape lpgroup}} by an L\texttt{\symbol{45}}presented or by a finitely presented group \mbox{\texttt{\mdseries\slshape H}}. The action of a generator of \mbox{\texttt{\mdseries\slshape H}} on \mbox{\texttt{\mdseries\slshape lpgroup}} is given by an automorphism in the list \mbox{\texttt{\mdseries\slshape auts}}. Thus for each generator of \mbox{\texttt{\mdseries\slshape H}} there must be an automorphism in the list \mbox{\texttt{\mdseries\slshape auts}}. 
\begin{Verbatim}[commandchars=!@|,fontsize=\small,frame=single,label=Example]
  !gapprompt@gap>| !gapinput@F := FreeGroup( "a" );|
  <free group on the generators [ a ]>
  !gapprompt@gap>| !gapinput@H := F / [ F.1^3 ];|
  <fp group on the generators [ a ]>
  !gapprompt@gap>| !gapinput@U := ExamplesOfLPresentations( 8 );|
  <L-presented group on the generators [ t, u, v ]>
  !gapprompt@gap>| !gapinput@aut:=GroupHomomorphismByImagesNC( U, U, [ U.1, U.2, U.3 ], [ U.2, U.3, U.1 ] );|
  [ t, u, v ] -> [ u, v, t ]
  !gapprompt@gap>| !gapinput@SplitExtensionByAutomorphismsLpGroup( U, H, [ aut ] );|
  <L-presented group on the generators [ t, u, v, a ]>
\end{Verbatim}
 }

 

\subsection{\textcolor{Chapter }{AsLpGroup}}
\logpage{[ 2, 6, 3 ]}\nobreak
\hyperdef{L}{X84F112247DA4037C}{}
{\noindent\textcolor{FuncColor}{$\triangleright$\enspace\texttt{AsLpGroup({\mdseries\slshape G})\index{AsLpGroup@\texttt{AsLpGroup}}
\label{AsLpGroup}
}\hfill{\scriptsize (operation)}}\\


 returns an ascending L\texttt{\symbol{45}}presentation for a finitely
presented group \mbox{\texttt{\mdseries\slshape G}} or for a free group \mbox{\texttt{\mdseries\slshape G}}. }

 

\subsection{\textcolor{Chapter }{IsomorphismLpGroup}}
\logpage{[ 2, 6, 4 ]}\nobreak
\hyperdef{L}{X856F237B7BAC3BC8}{}
{\noindent\textcolor{FuncColor}{$\triangleright$\enspace\texttt{IsomorphismLpGroup({\mdseries\slshape G})\index{IsomorphismLpGroup@\texttt{IsomorphismLpGroup}}
\label{IsomorphismLpGroup}
}\hfill{\scriptsize (operation)}}\\


 returns an isomorphism from a finitely presented group \mbox{\texttt{\mdseries\slshape G}} or from a free group \mbox{\texttt{\mdseries\slshape G}} to the L\texttt{\symbol{45}}presented group obtained from the method \texttt{AsLpGroup} (\ref{AsLpGroup}). 
\begin{Verbatim}[commandchars=!@|,fontsize=\small,frame=single,label=Example]
  !gapprompt@gap>| !gapinput@F:=FreeGroup( 2 );|
  <free group on the generators [ f1, f2 ]>
  !gapprompt@gap>| !gapinput@G:=F/[ F.1^2, F.2^2, Comm( F.1, F.2 ) ];|
  <fp group on the generators [ f1, f2 ]>
  !gapprompt@gap>| !gapinput@IsomorphismLpGroup( G );|
  [ f1, f2 ] -> [ f1, f2 ]
  !gapprompt@gap>| !gapinput@Range(last);|
  <L-presented group on the generators [ f1, f2 ]>
  !gapprompt@gap>| !gapinput@Display(last);|
  generators = [ f1, f2 ]
  fixed relators = [ ]
  endomorphism = [
  IdentityMapping( <free group on the generators [ f1, f2 ]> ) ]
  iterated relators = [
  f1^2,
  f2^2,
  f1^-1*f2^-1*f1*f2 ]
\end{Verbatim}
 }

 }

 }

 
\chapter{\textcolor{Chapter }{Nilpotent Quotients of L\texttt{\symbol{45}}presented groups}}\logpage{[ 3, 0, 0 ]}
\hyperdef{L}{X824CC9CA824D3F1E}{}
{
 Our nilpotent quotient algorithm for finitely L\texttt{\symbol{45}}presented
groups generalizes Nickel's algorithm for finitely presented groups; see \cite{Nickel96}. It determines a nilpotent presentation for the lower central series quotient
of an invariantly L\texttt{\symbol{45}}presented group. A nilpotent
presentation is a polycyclic presentation whose polycyclic series refines the
lower central series of the group (see the description in the \textsf{NQ}\texttt{\symbol{45}}package for further details). In general, our algorithm
determines a polycyclic presentation for the nilpotent quotient of an
arbitrary finitely L\texttt{\symbol{45}}presented group. For further details
on our algorithm we refer to \cite{BEH08} or to the diploma thesis \cite{H08}. 
\section{\textcolor{Chapter }{New methods for L\texttt{\symbol{45}}presented groups}}\logpage{[ 3, 1, 0 ]}
\hyperdef{L}{X791C3E5280F38329}{}
{
 

\subsection{\textcolor{Chapter }{NilpotentQuotient}}
\logpage{[ 3, 1, 1 ]}\nobreak
\hyperdef{L}{X8216791583DE512C}{}
{\noindent\textcolor{FuncColor}{$\triangleright$\enspace\texttt{NilpotentQuotient({\mdseries\slshape g[, c]})\index{NilpotentQuotient@\texttt{NilpotentQuotient}}
\label{NilpotentQuotient}
}\hfill{\scriptsize (operation)}}\\


 returns a polycyclic presentation for the class\texttt{\symbol{45}}\mbox{\texttt{\mdseries\slshape c}} quotient $g/\gamma_{c+1}(g)$ of the L\texttt{\symbol{45}}presented group \mbox{\texttt{\mdseries\slshape g}} if \mbox{\texttt{\mdseries\slshape c}} is specified. If \mbox{\texttt{\mdseries\slshape c}} is not given, this method attempts to compute the largest nilpotent quotient
of \mbox{\texttt{\mdseries\slshape g}} and will terminate only if \mbox{\texttt{\mdseries\slshape g}} has a largest nilpotent quotient.

 The following example computes the class\texttt{\symbol{45}}5 quotient of the
Grigorchuk group. 
\begin{Verbatim}[commandchars=!@|,fontsize=\small,frame=single,label=Example]
  !gapprompt@gap>| !gapinput@G := ExamplesOfLPresentations( 1 );;|
  !gapprompt@gap>| !gapinput@H := NilpotentQuotient( G, 5 );|
  Pcp-group with orders [ 2, 2, 2, 2, 2, 2, 2, 2, 2, 2 ] 
  !gapprompt@gap>| !gapinput@lcs := LowerCentralSeries( H );|
  [ Pcp-group with orders [ 2, 2, 2, 2, 2, 2, 2, 2, 2, 2 ], 
    Pcp-group with orders [ 2, 2, 2, 2, 2, 2, 2 ], 
    Pcp-group with orders [ 2, 2, 2, 2, 2 ], Pcp-group with orders [ 2, 2, 2 ], 
    Pcp-group with orders [ 2, 2 ], Pcp-group with orders [  ] ]
  !gapprompt@gap>| !gapinput@List( [ 1..5 ], x -> lcs[ x ] / lcs[ x+1 ] ); |
  [ Pcp-group with orders [ 2, 2, 2 ], Pcp-group with orders [ 2, 2 ], 
    Pcp-group with orders [ 2, 2 ], Pcp-group with orders [ 2 ], 
    Pcp-group with orders [ 2, 2 ] ]
\end{Verbatim}
 }

 

\subsection{\textcolor{Chapter }{LargestNilpotentQuotient}}
\logpage{[ 3, 1, 2 ]}\nobreak
\hyperdef{L}{X79AC8BE285CBB392}{}
{\noindent\textcolor{FuncColor}{$\triangleright$\enspace\texttt{LargestNilpotentQuotient({\mdseries\slshape g})\index{LargestNilpotentQuotient@\texttt{LargestNilpotentQuotient}}
\label{LargestNilpotentQuotient}
}\hfill{\scriptsize (operation)}}\\


 returns the largest nilpotent quotient of the L\texttt{\symbol{45}}presented
group \mbox{\texttt{\mdseries\slshape g}} if it exists. It uses the method \texttt{NilpotentQuotient} (\ref{NilpotentQuotient}). If \mbox{\texttt{\mdseries\slshape g}} has no largest nilpotent quotient, this method will not terminate. }

 

\subsection{\textcolor{Chapter }{NqEpimorphismNilpotentQuotient}}
\logpage{[ 3, 1, 3 ]}\nobreak
\hyperdef{L}{X8758F663782AE655}{}
{\noindent\textcolor{FuncColor}{$\triangleright$\enspace\texttt{NqEpimorphismNilpotentQuotient({\mdseries\slshape g[, p/c]})\index{NqEpimorphismNilpotentQuotient@\texttt{NqEpimorphismNilpotentQuotient}}
\label{NqEpimorphismNilpotentQuotient}
}\hfill{\scriptsize (operation)}}\\


 This method returns an epimorphism from the L\texttt{\symbol{45}}presented
group \mbox{\texttt{\mdseries\slshape g}} onto a nilpotent quotient. If the optional argument is an integer \mbox{\texttt{\mdseries\slshape c}}, the epimorphism is onto the maximal class\texttt{\symbol{45}}\mbox{\texttt{\mdseries\slshape c}} quotient $g//\gamma_{c+1}(g)$. If no second argument is given, this method attempts to compute an
epimorphism onto the largest nilpotent quotient of \mbox{\texttt{\mdseries\slshape g}}. If \mbox{\texttt{\mdseries\slshape g}} does not have a largest nilpotent quotient, this method will not terminate.

 If a pcp\texttt{\symbol{45}}group \mbox{\texttt{\mdseries\slshape p}} is given as additional parameter, then \mbox{\texttt{\mdseries\slshape p}} has to be a nilpotent quotient of \mbox{\texttt{\mdseries\slshape g}}. The method computes an epimorphism from the L\texttt{\symbol{45}}presented
group \mbox{\texttt{\mdseries\slshape g}} onto \mbox{\texttt{\mdseries\slshape p}}.

 The following example computes an epimorphism from the Grigorchuk group onto
its class\texttt{\symbol{45}}5, class\texttt{\symbol{45}}7, and
class\texttt{\symbol{45}}10 quotients. 
\begin{Verbatim}[commandchars=!@|,fontsize=\small,frame=single,label=Example]
  !gapprompt@gap>| !gapinput@G := ExamplesOfLPresentations( 1 );|
  <L-presented group on the generators [ a, b, c, d ]>
  !gapprompt@gap>| !gapinput@epi := NqEpimorphismNilpotentQuotient( G, 5 );|
  [ a, b, c, d ] -> [ g1, g2*g3, g2, g3 ]
  !gapprompt@gap>| !gapinput@H := Image( epi );|
  Pcp-group with orders [ 2, 2, 2, 2, 2, 2, 2, 2, 2, 2 ]
  !gapprompt@gap>| !gapinput@NilpotencyClassOfGroup( H );|
  5
  !gapprompt@gap>| !gapinput@H := NilpotentQuotient( G, 7 );|
  Pcp-group with orders [ 2, 2, 2, 2, 2, 2, 2, 2, 2, 2, 2, 2, 2 ]
  !gapprompt@gap>| !gapinput@NilpotentQuotient( G, 10 );|
  Pcp-group with orders [ 2, 2, 2, 2, 2, 2, 2, 2, 2, 2, 2, 2, 2, 2, 2, 2, 2, 2 ]
  !gapprompt@gap>| !gapinput@NqEpimorphismNilpotentQuotient( G, H );|
  [ a, b, c, d ] -> [ g1, g2*g3, g2, g3 ]
  !gapprompt@gap>| !gapinput@Image( last );|
  Pcp-group with orders [ 2, 2, 2, 2, 2, 2, 2, 2, 2, 2, 2, 2, 2 ]
\end{Verbatim}
 }

 

\subsection{\textcolor{Chapter }{AbelianInvariants}}
\logpage{[ 3, 1, 4 ]}\nobreak
\hyperdef{L}{X812827937F403300}{}
{\noindent\textcolor{FuncColor}{$\triangleright$\enspace\texttt{AbelianInvariants({\mdseries\slshape g})\index{AbelianInvariants@\texttt{AbelianInvariants}}
\label{AbelianInvariants}
}\hfill{\scriptsize (operation)}}\\


 computes the abelian invariants of the L\texttt{\symbol{45}}presented group \mbox{\texttt{\mdseries\slshape g}}. It uses the operation \texttt{NilpotentQuotient} (\ref{NilpotentQuotient}). 
\begin{Verbatim}[commandchars=!@|,fontsize=\small,frame=single,label=Example]
  !gapprompt@gap>| !gapinput@G := ExamplesOfLPresentations( 1 );;|
  !gapprompt@gap>| !gapinput@AbelianInvariants( G );|
  [ 2, 2, 2 ]
\end{Verbatim}
 }

 }

 
\section{\textcolor{Chapter }{A brief description of the algorithm}}\logpage{[ 3, 2, 0 ]}
\hyperdef{L}{X7C529DA9802E603E}{}
{
 In the following we give a brief description of the nilpotent quotient
algorithm for an arbitrary finitely L\texttt{\symbol{45}}presented group. For
further details, we refer to \cite{BEH08} and the diploma thesis \cite{H08}.

 Let $(S,Q,\Phi,R)$ be a finite L\texttt{\symbol{45}}presentation defining the
L\texttt{\symbol{45}}presented group $G$ and let $(S,Q',\Phi,R)$ be an underlying invariant L\texttt{\symbol{45}}presentation. Write $\bar G$ for the invariantly L\texttt{\symbol{45}}presented group defined by $(S,Q',\Phi,R)$.

 The first step in computing a polycyclic presentation for $G/\gamma_{c+1}(G)$ is to determine a nilpotent presentation for $\bar G/\gamma_{c+1}(\bar G)$. This will be done by induction on $c$. The induction step of our algorithm generalizes the induction step of
Nickel's algorithm which mainly relies on Hermite normal form computations. In
order to use this rather fast linear algebra, we must require the group to be
invariantly L\texttt{\symbol{45}}presented. Therefore, the fixed relators must
be handled separately by reducing to an underlying invariant
L\texttt{\symbol{45}}presentation first.

 The induction step of our algorithm then returns a nilpotent presentation $H$ for the quotient $\bar G/\gamma_{c+1}(\bar G)$ and an epimorphism $\delta\colon\bar G\to H$. Both are used to determine a polycyclic presentation for the nilpotent
quotient $G/\gamma_{c+1}(G)$ using an extension $\delta'\colon F_S\to H$ of the epimorphism $\delta$. The quotient $G/\gamma_{c+1}(G)$ is isomorphic to the factor group $H/\langle Q^{\delta'}\rangle^H$. We use the \textsf{Polycyclic}\texttt{\symbol{45}}package to compute a polycyclic presentation for $H/\langle Q^{\delta'}\rangle^H$.

 The efficiency of this general approach depends on the underlying invariant
L\texttt{\symbol{45}}presentation $(S,Q',\Phi,R)$. The set of fixed relators $Q'$ should be as large as possible. Otherwise, the nilpotent quotient $H$ can be large even if the nilpotent quotient $G/\gamma_{c+1}(G)$ is rather small.

 The following example demonstrates the different behavior of our nilpotent
quotient algorithm for the Grigorchuk group with its finite
L\texttt{\symbol{45}}presentation 
\[\Big(\{a,c,b,d\},\{a^2,b^2,c^2,d^2,bcd\},\{\sigma\},\{[d,d^a],[d,d^{acaca}]\}
\Big). \]
 This latter L\texttt{\symbol{45}}presentation is obviously an invariant
L\texttt{\symbol{45}}presentation. Hence, we can either use the property \texttt{IsInvariantLPresentation} (\ref{IsInvariantLPresentation}) or the attribute \texttt{UnderlyingInvariantLPresentation} (\ref{UnderlyingInvariantLPresentation}). First, one has to construct the group as described in Section \ref{createlp}: 
\begin{Verbatim}[commandchars=!@|,fontsize=\small,frame=single,label=Example]
  !gapprompt@gap>| !gapinput@F := FreeGroup( "a", "b", "c", "d" );|
  <free group on the generators [ a, b, c, d ]>
  !gapprompt@gap>| !gapinput@AssignGeneratorVariables( F );|
  #I  Assigned the global variables [ a, b, c, d ]
  !gapprompt@gap>| !gapinput@rels := [ a^2, b^2, c^2, d^2, b*d*c ];;|
  !gapprompt@gap>| !gapinput@endos := [ GroupHomomorphismByImagesNC( F, F, [ a, b, c, d ], [ c^a, d, b, c ]) ];;|
  !gapprompt@gap>| !gapinput@itrels := [ Comm( d, d^a ), Comm( d, d^(a*c*a*c*a) ) ];;|
  !gapprompt@gap>| !gapinput@G := LPresentedGroup( F, rels, endos, itrels );|
  <L-presented group on the generators [ a, b, c, d ]>
  !gapprompt@gap>| !gapinput@List( rels, x -> x^endos[1] );|
  [ a^-1*c^2*a, d^2, b^2, c^2, d*c*b ]
\end{Verbatim}
 The property \texttt{IsInvariantLPresentation} (\ref{IsInvariantLPresentation}) can be set manually using \texttt{SetInvariantLPresentation}: 
\begin{Verbatim}[commandchars=!@|,fontsize=\small,frame=single,label=Example]
  !gapprompt@gap>| !gapinput@SetIsInvariantLPresentation( G, true );|
  !gapprompt@gap>| !gapinput@NilpotentQuotient( G, 4 );|
  Pcp-group with orders [ 2, 2, 2, 2, 2, 2, 2, 2 ]
  !gapprompt@gap>| !gapinput@StringTime( time );|
  " 0:00:00.032"
\end{Verbatim}
 On the other hand, one can use the attribute \texttt{UnderlyingInvariantLPresentation} (\ref{UnderlyingInvariantLPresentation}) as follows: 
\begin{Verbatim}[commandchars=!@|,fontsize=\small,frame=single,label=Example]
  !gapprompt@gap>| !gapinput@U := LPresentedGroup( F, rels, endos, itrels );|
  <L-presented group on the generators [ a, b, c, d ]>
  !gapprompt@gap>| !gapinput@SetUnderlyingInvariantLPresentation( G, U );|
  !gapprompt@gap>| !gapinput@NilpotentQuotient( G, 4 );|
  Pcp-group with orders [ 2, 2, 2, 2, 2, 2, 2, 2 ]
  !gapprompt@gap>| !gapinput@StringTime( time );|
  " 0:00:00.028"
\end{Verbatim}
 For saving memory the first method should be preferred in this case. In
general, the L\texttt{\symbol{45}}presentation is not invariant (or not known
to be invariant) and thus the underlying invariant
L\texttt{\symbol{45}}presentation has fewer fixed relators than the group $G$ itself. In this case, the second method is the method of choice.

 There is a brute\texttt{\symbol{45}}force method implemented for the operation \texttt{UnderlyingInvariantLPresentation} (\ref{UnderlyingInvariantLPresentation}) which works quite well on the \texttt{ExamplesOfLPresentations} (\ref{ExamplesOfLPresentations}). However, in the worst case, this method will return the underlying ascending
L\texttt{\symbol{45}}presentation. The following example shows the influence
of this choice to the runtime of the nilpotent quotient algorithm. After
defining the group $G$ as above, we set the attribute \texttt{UnderlyingInvariantLPresentation} (\ref{UnderlyingInvariantLPresentation}) as follows: 
\begin{Verbatim}[commandchars=!@|,fontsize=\small,frame=single,label=Example]
  !gapprompt@gap>| !gapinput@SetUnderlyingInvariantLPresentation( G, UnderlyingAscendingLPresentation(G) );|
  !gapprompt@gap>| !gapinput@NilpotentQuotient( G, 4 );|
  Pcp-group with orders [ 2, 2, 2, 2, 2, 2, 2, 2 ]
  !gapprompt@gap>| !gapinput@StringTime( time );|
  " 0:00:02.700"
\end{Verbatim}
 }

 
\section{\textcolor{Chapter }{Nilpotent Quotient Systems for invariant L\texttt{\symbol{45}}presentations}}\logpage{[ 3, 3, 0 ]}
\hyperdef{L}{X864A3F6F796E99DF}{}
{
 For an invariantly L\texttt{\symbol{45}}presented group $G$, our algorithm computes a nilpotent presentation for $G/\gamma_{c+1}(G)$ by computing a \mbox{\texttt{\mdseries\slshape weighted nilpotent quotient system}} for $G/G'$ and extending it inductively to a weighted nilpotent quotient system for $G/\gamma_{c+1}(G)$.

 In the \textsf{lpres} package, a weighted nilpotent quotient system is a record containing the
following entries: 
\begin{description}
\item[{Lpres}] the invariantly L\texttt{\symbol{45}}presented group $G$.
\item[{Pccol}] \texttt{FromTheLeftCollector} (\textbf{polycyclic: FromTheLeftCollector}) of the nilpotent quotient represented by this quotient system.
\item[{Imgs}] the images of the generators of the L\texttt{\symbol{45}}presented group $G$ under the epimorphism onto the nilpotent quotient \mbox{\texttt{\mdseries\slshape Pccol}}. For each generator of $G$ there is an integer or a generator exponent list. If the image is an integer \mbox{\texttt{\mdseries\slshape int}}, the image is a definition of the \mbox{\texttt{\mdseries\slshape int}}\texttt{\symbol{45}}th generator of the nilpotent presentation \mbox{\texttt{\mdseries\slshape Pccol}}.
\item[{Epimorphism}] an epimorphism from the L\texttt{\symbol{45}}presented group $G$ onto its nilpotent quotient \mbox{\texttt{\mdseries\slshape Pccol}} with the images of the generators given by \mbox{\texttt{\mdseries\slshape Imgs}}.
\item[{Weights}] a list of the weight of each generator of the nilpotent presentation \mbox{\texttt{\mdseries\slshape Pccol}}.
\item[{Definitions}] the definition of each generator of \mbox{\texttt{\mdseries\slshape Pccol}}. Each generator in the quotient system has a definition as an image or as a
commutator of the form $[a_j,a_i]$ where $a_j$ and $a_i$ are generators of a certain weight. If the \mbox{\texttt{\mdseries\slshape i}}\texttt{\symbol{45}}th entry is an integer, the \mbox{\texttt{\mdseries\slshape i}}\texttt{\symbol{45}}th generator of \mbox{\texttt{\mdseries\slshape Pccol}} has a definition as an image. Otherwise, the \mbox{\texttt{\mdseries\slshape i}}\texttt{\symbol{45}}th entry is a $2$\texttt{\symbol{45}}tuple $[k,l]$ and the \mbox{\texttt{\mdseries\slshape i}}\texttt{\symbol{45}}th generator has a definition as commutator $[a_k,a_l]$.
\end{description}
 A weighted nilpotent quotient system of an invariantly
L\texttt{\symbol{45}}presented group can be computed with the following
functions. 

\subsection{\textcolor{Chapter }{InitQuotientSystem}}
\logpage{[ 3, 3, 1 ]}\nobreak
\hyperdef{L}{X7E58D47A8729FA8E}{}
{\noindent\textcolor{FuncColor}{$\triangleright$\enspace\texttt{InitQuotientSystem({\mdseries\slshape lpgroup})\index{InitQuotientSystem@\texttt{InitQuotientSystem}}
\label{InitQuotientSystem}
}\hfill{\scriptsize (operation)}}\\


 computes a weighted nilpotent quotient system for the abelian quotient of the
L\texttt{\symbol{45}}presented group \mbox{\texttt{\mdseries\slshape lpgroup}}. }

 

\subsection{\textcolor{Chapter }{ExtendQuotientSystem}}
\logpage{[ 3, 3, 2 ]}\nobreak
\hyperdef{L}{X7910D0698781E02A}{}
{\noindent\textcolor{FuncColor}{$\triangleright$\enspace\texttt{ExtendQuotientSystem({\mdseries\slshape QS})\index{ExtendQuotientSystem@\texttt{ExtendQuotientSystem}}
\label{ExtendQuotientSystem}
}\hfill{\scriptsize (operation)}}\\


 extends the weighted nilpotent quotient system \mbox{\texttt{\mdseries\slshape QS}} for a class\texttt{\symbol{45}}$c$ quotient of an invariantly L\texttt{\symbol{45}}presented group to a weighted
nilpotent quotient system of its class\texttt{\symbol{45}}$c+1$ quotient. 
\begin{Verbatim}[commandchars=!@|,fontsize=\small,frame=single,label=Example]
  !gapprompt@gap>| !gapinput@G := ExamplesOfLPresentations( 1 );|
  <L-presented group on the generators [ a, b, c, d ]>
  !gapprompt@gap>| !gapinput@Q := InitQuotientSystem( G );|
  rec( Lpres := <L-presented group on the generators [ a, b, c, d ]>, 
    Pccol := <<from the left collector with 3 generators>>, 
    Imgs := [ 1, [ 2, 1, 3, 1 ], 2, 3 ], Epimorphism := [ a, b, c, d ] -> 
      [ g1, g2*g3, g2, g3 ], Weights := [ 1, 1, 1 ], Definitions := [ 1, 3, 4 ] 
   )
  !gapprompt@gap>| !gapinput@ExtendQuotientSystem( Q );|
  rec( Lpres := <L-presented group on the generators [ a, b, c, d ]>, 
    Pccol := <<from the left collector with 5 generators>>, 
    Imgs := [ 1, [ 2, 1, 3, 1 ], 2, 3 ], 
    Definitions := [ 1, 3, 4, [ 2, 1 ], [ 3, 1 ] ], 
    Weights := [ 1, 1, 1, 2, 2 ], Epimorphism := [ a, b, c, d ] -> 
      [ g1, g2*g3, g2, g3 ] )
\end{Verbatim}
 }

 }

 
\section{\textcolor{Chapter }{Attributes of L\texttt{\symbol{45}}presented groups related with the nilpotent
quotient algorithm}}\logpage{[ 3, 4, 0 ]}
\hyperdef{L}{X87CA2F188762A2B5}{}
{


 To avoid repeated extensions of a weighted nilpotent quotient system the
largest known quotient system is stored as an attribute of the invariantly
L\texttt{\symbol{45}}presented group. For non\texttt{\symbol{45}}invariantly
L\texttt{\symbol{45}}presented groups (or groups which are not known to be
invariantly L\texttt{\symbol{45}}presented) the known epimorphisms onto the
nilpotent quotients are stored as an attribute. 

\subsection{\textcolor{Chapter }{NilpotentQuotientSystem}}
\logpage{[ 3, 4, 1 ]}\nobreak
\hyperdef{L}{X7CC4586B85C22457}{}
{\noindent\textcolor{FuncColor}{$\triangleright$\enspace\texttt{NilpotentQuotientSystem({\mdseries\slshape lpgroup})\index{NilpotentQuotientSystem@\texttt{NilpotentQuotientSystem}}
\label{NilpotentQuotientSystem}
}\hfill{\scriptsize (attribute)}}\\


 stores the largest known weighted nilpotent quotient system of an invariantly
L\texttt{\symbol{45}}presented group. 
\begin{Verbatim}[commandchars=!@|,fontsize=\small,frame=single,label=Example]
  !gapprompt@gap>| !gapinput@G := ExamplesOfLPresentations( 1 );;|
  !gapprompt@gap>| !gapinput@NilpotentQuotient( G, 5 );|
  Pcp-group with orders [ 2, 2, 2, 2, 2, 2, 2, 2, 2, 2 ]
  !gapprompt@gap>| !gapinput@NilpotentQuotientSystem( G );|
  rec( Lpres := <L-presented group on the generators [ a, b, c, d ]>, 
    Pccol := <<from the left collector with 10 generators>>, 
    Imgs := [ 1, [ 2, 1, 3, 1 ], 2, 3 ], 
    Definitions := [ 1, 3, 4, [ 2, 1 ], [ 3, 1 ], [ 4, 2 ], [ 4, 3 ], [ 7, 1 ], 
        [ 8, 2 ], [ 8, 3 ] ], Weights := [ 1, 1, 1, 2, 2, 3, 3, 4, 5, 5 ], 
    Epimorphism := [ a, b, c, d ] -> [ g1, g2*g3, g2, g3 ] )
  !gapprompt@gap>| !gapinput@NilpotencyClassOfGroup( PcpGroupByCollectorNC( last.Pccol ) );|
  5
\end{Verbatim}
 }

 

\subsection{\textcolor{Chapter }{NilpotentQuotients}}
\logpage{[ 3, 4, 2 ]}\nobreak
\hyperdef{L}{X7D54126783CB7118}{}
{\noindent\textcolor{FuncColor}{$\triangleright$\enspace\texttt{NilpotentQuotients({\mdseries\slshape lpgroup})\index{NilpotentQuotients@\texttt{NilpotentQuotients}}
\label{NilpotentQuotients}
}\hfill{\scriptsize (attribute)}}\\


 stores all known epimorphisms onto the nilpotent quotients of \mbox{\texttt{\mdseries\slshape lpgroup}}. The nilpotent quotients are accessible by the operation \texttt{Range} (\textbf{Reference: Range of a general mapping}). 
\begin{Verbatim}[commandchars=!@|,fontsize=\small,frame=single,label=Example]
  !gapprompt@gap>| !gapinput@G:=ExamplesOfLPresentations( 3 );;|
  !gapprompt@gap>| !gapinput@HasIsInvariantLPresentation( G );|
  false
  !gapprompt@gap>| !gapinput@NilpotentQuotient( G, 3 );|
  Pcp-group with orders [ 0, 2, 2, 2 ]
  !gapprompt@gap>| !gapinput@NilpotentQuotients( G );|
  [ [ a, t, u ] -> [ g2, g1, g2 ], [ a, t, u ] -> [ g2, g1, g2 ],
    [ a, t, u ] -> [ g2, g1, g2 ] ]
  !gapprompt@gap>| !gapinput@Range( last[2] );|
  Pcp-group with orders [ 0, 2, 2 ]
\end{Verbatim}
 The underlying invariant L\texttt{\symbol{45}}presentation has stored its
largest weighted nilpotent quotient system as an attribute. 
\begin{Verbatim}[commandchars=!@|,fontsize=\small,frame=single,label=Example]
  !gapprompt@gap>| !gapinput@NilpotentQuotientSystem( UnderlyingInvariantLPresentation( G ) );|
  rec( Lpres := <L-presented group on the generators [ a, t, u ]>,
    Pccol := <<from the left collector with 9 generators>>, Imgs := [ 1, 2, 3 ],
    Definitions := [ 1, 2, 3, [ 2, 1 ], [ 3, 2 ], [ 4, 1 ], [ 4, 2 ], [ 5, 2 ],
        [ 5, 3 ] ], Weights := [ 1, 1, 1, 2, 2, 3, 3, 3, 3 ],
    Epimorphism := [ a, t, u ] -> [ g1, g2, g3 ] )
\end{Verbatim}
 }

 }

 
\section{\textcolor{Chapter }{The Info\texttt{\symbol{45}}Class InfoLPRES}}\logpage{[ 3, 5, 0 ]}
\hyperdef{L}{X7BB56B4C7C1EFAB8}{}
{


 To get some information about the progress of the algorithm, one can use the
info class \texttt{InfoLPRES} (\ref{InfoLPRES}). 

\subsection{\textcolor{Chapter }{InfoLPRES}}
\logpage{[ 3, 5, 1 ]}\nobreak
\hyperdef{L}{X85F6BC1F8573D710}{}
{\noindent\textcolor{FuncColor}{$\triangleright$\enspace\texttt{InfoLPRES\index{InfoLPRES@\texttt{InfoLPRES}}
\label{InfoLPRES}
}\hfill{\scriptsize (info class)}}\\


 is the info class of the \textsf{lpres}\texttt{\symbol{45}}package. If the info\texttt{\symbol{45}}level is $1$, the info\texttt{\symbol{45}}class gives further information on the progress
of the nilpotent quotient algorithm for L\texttt{\symbol{45}}presented groups.
The info\texttt{\symbol{45}}level $2$ also includes some information on the runtime of our algorithm while the
info\texttt{\symbol{45}}level $3$ is mainly used for debugging\texttt{\symbol{45}}purposes. An example of such a
session for the Grigorchuk group is shown below: 
\begin{Verbatim}[commandchars=!@|,fontsize=\small,frame=single,label=Example]
  !gapprompt@gap>| !gapinput@SetInfoLevel( InfoLPRES, 1 );;|
  !gapprompt@gap>| !gapinput@G:=ExamplesOfLPresentations( 1 );|
  #I  The Grigorchuk group on 4 generators
  <L-presented group on the generators [ a, b, c, d ]>
  !gapprompt@gap>| !gapinput@NilpotentQuotient( G, 3 );|
  #I  Class 1: 3 generators with relative orders: [ 2, 2, 2 ]
  #I  Class 2: 2 generators with relative orders: [ 2, 2 ]
  #I  Class 3: 2 generators with relative orders: [ 2, 2 ]
  Pcp-group with orders [ 2, 2, 2, 2, 2, 2, 2 ]
  !gapprompt@gap>| !gapinput@SetInfoLevel( InfoLPRES, 2 );|
  !gapprompt@gap>| !gapinput@NilpotentQuotient( G, 5 );|
  #I  Time spent for spinning algo:  0:00:00.004
  #I  Class 4: 1 generators with relative orders: [ 2 ]
  #I  Runtime for this step  0:00:00.028
  #I  Time spent for spinning algo:  0:00:00.008
  #I  Class 5: 2 generators with relative orders: [ 2, 2 ]
  #I  Runtime for this step  0:00:00.036
  Pcp-group with orders [ 2, 2, 2, 2, 2, 2, 2, 2, 2, 2 ]
\end{Verbatim}
 }

 

\subsection{\textcolor{Chapter }{InfoLPRES{\textunderscore}MAX{\textunderscore}GENS}}
\logpage{[ 3, 5, 2 ]}\nobreak
\hyperdef{L}{X80F8139B81D2294E}{}
{\noindent\textcolor{FuncColor}{$\triangleright$\enspace\texttt{InfoLPRES{\textunderscore}MAX{\textunderscore}GENS\index{InfoLPRES{\textunderscore}MAX{\textunderscore}GENS@\texttt{Info}\-\texttt{L}\-\texttt{P}\-\texttt{R}\-\texttt{E}\-\texttt{S{\textunderscore}}\-\texttt{M}\-\texttt{A}\-\texttt{X{\textunderscore}}\-\texttt{GENS}}
\label{InfoLPRESuScoreMAXuScoreGENS}
}\hfill{\scriptsize (global variable)}}\\


 this global variable sets the limit of generators whose relative order will be
shown on each step of the nilpotent quotient algorithm, if the
info\texttt{\symbol{45}}level of \texttt{InfoLPRES} (\ref{InfoLPRES}) is positive. }

 }

 }

 
\chapter{\textcolor{Chapter }{Subgroups of L\texttt{\symbol{45}}presented groups}}\logpage{[ 4, 0, 0 ]}
\hyperdef{L}{X874D64AA789F224E}{}
{


 As shown in \cite{MR2864562} it is possible to deal with finite index subgroups of
L\texttt{\symbol{45}}presented groups algorithmically. The \textsf{lpres}\texttt{\symbol{45}}package provides straightforward methods to deal with
these subgroups.

 The Reidemeister\texttt{\symbol{45}}Schreier algorithm from \cite{MR2876891} is implemented, and computes presentations for such subgroups. 
\section{\textcolor{Chapter }{Creating a subgroup of an L\texttt{\symbol{45}}presented group}}\logpage{[ 4, 1, 0 ]}
\hyperdef{L}{X86B9E4BD7F5D1610}{}
{
 There are two ways of defining subgroups of finite index of an \mbox{\texttt{\mdseries\slshape lpgroup}}. The first is to define the subgroup by its generators while the second
defines the subgroup by a coset\texttt{\symbol{45}}table. Generators of
subgroup of the latter type can be computed with the usual
Schreier\texttt{\symbol{45}}algorithm. 

\subsection{\textcolor{Chapter }{Subgroup}}
\logpage{[ 4, 1, 1 ]}\nobreak
\hyperdef{L}{X7C82AA387A42DCA0}{}
{\noindent\textcolor{FuncColor}{$\triangleright$\enspace\texttt{Subgroup({\mdseries\slshape G, gens})\index{Subgroup@\texttt{Subgroup}}
\label{Subgroup}
}\hfill{\scriptsize (function)}}\\


 creates the subgroup \mbox{\texttt{\mdseries\slshape U}} of \mbox{\texttt{\mdseries\slshape G}} generated by \mbox{\texttt{\mdseries\slshape gens}}. The Parent value of \mbox{\texttt{\mdseries\slshape U}} will be \mbox{\texttt{\mdseries\slshape G}}.

 For example, the branching subgroup of the Grigorchuk group can be defined as
follows, and a presentation can be computed using \texttt{IsomorphismLpGroup} (\ref{IsomorphismLpGroup}): 
\begin{Verbatim}[commandchars=!@|,fontsize=\small,frame=single,label=Example]
  !gapprompt@gap>| !gapinput@G := ExamplesOfLPresentations(1);;|
  !gapprompt@gap>| !gapinput@a := G.1;; b := G.2;; c := G.3;; d := G.4;;|
  !gapprompt@gap>| !gapinput@K := Subgroup( G, [ Comm( a, b ), Comm( b^a, d ), Comm( b, d^a ) ] );|
  Group([ a^-1*b^-1*a*b, b^-1*a^-1*d^-1*a*b*a^-1*d*a, a^-1*b^-1*a*d^-1*a^-1*b*a*d ])
  !gapprompt@gap>| !gapinput@iso := IsomorphismLpGroup(K);|
  [ a^-1*b^-1*a*b, a^-1*b^-1*a*d^-1*a^-1*b*a*d, b^-1*a^-1*d^-1*a*b*a^-1*d*a ] ->
  [ x1^-1*x13^-1*x12*x3, x1^-1*x13^-1*x12*x8^-1*x22^-1*x23*x24*x11, x3^-1*x12^-1*x18^-1*x29*x21*x1 ]
  !gapprompt@gap>| !gapinput@Range(iso);|
  <LpGroup with 98 generators>
\end{Verbatim}
 }

 

\subsection{\textcolor{Chapter }{SubgroupLpGroupByCosetTable}}
\logpage{[ 4, 1, 2 ]}\nobreak
\hyperdef{L}{X7FC6C908782DEA48}{}
{\noindent\textcolor{FuncColor}{$\triangleright$\enspace\texttt{SubgroupLpGroupByCosetTable({\mdseries\slshape G, Tab})\index{SubgroupLpGroupByCosetTable@\texttt{SubgroupLpGroupByCosetTable}}
\label{SubgroupLpGroupByCosetTable}
}\hfill{\scriptsize (operation)}}\\


 creates the subgroup \mbox{\texttt{\mdseries\slshape U}} of \mbox{\texttt{\mdseries\slshape G}} which is represented by the coset\texttt{\symbol{45}}table \mbox{\texttt{\mdseries\slshape Tab}}.

 For instance, the branching subgroup of the Grigorchuk group can be defined by
the following coset\texttt{\symbol{45}}table: 
\begin{Verbatim}[commandchars=!@|,fontsize=\small,frame=single,label=Example]
  !gapprompt@gap>| !gapinput@Tab := [ [ 2, 1, 6, 9, 10, 3, 11, 12, 4, 5, 7, 8, 15, 16, 13, 14 ],|
    [ 2, 1, 6, 9, 10, 3, 11, 12, 4, 5, 7, 8, 15, 16, 13, 14 ],
    [ 3, 6, 1, 5, 4, 2, 8, 7, 10, 9, 12, 11, 14, 13, 16, 15 ],
    [ 3, 6, 1, 5, 4, 2, 8, 7, 10, 9, 12, 11, 14, 13, 16, 15 ],
    [ 4, 7, 5, 1, 3, 8, 2, 6, 13, 14, 15, 16, 9, 10, 11, 12 ],
    [ 4, 7, 5, 1, 3, 8, 2, 6, 13, 14, 15, 16, 9, 10, 11, 12 ],
    [ 5, 8, 4, 3, 1, 7, 6, 2, 14, 13, 16, 15, 10, 9, 12, 11 ],
    [ 5, 8, 4, 3, 1, 7, 6, 2, 14, 13, 16, 15, 10, 9, 12, 11 ] ]
  !gapprompt@gap>| !gapinput@U := SubgroupLpGroupByCosetTable( G, Tab );|
  Group(<A>subgroup of L-presented group, no generators known</A>)
  !gapprompt@gap>| !gapinput@U = K;|
  true
\end{Verbatim}
 The generators of \mbox{\texttt{\mdseries\slshape U}} can be computed with the Schreier\texttt{\symbol{45}}algorithm which is
implemented in the method \texttt{GeneratorsOfGroup} (\textbf{Reference: GeneratorsOfGroup}). }

 }

 
\section{\textcolor{Chapter }{Computing the index of finite\texttt{\symbol{45}}index subgroups}}\logpage{[ 4, 2, 0 ]}
\hyperdef{L}{X7A4EB4E0819ACB91}{}
{


 In principle, it is possible to compute the index of a finite index subgroup
of an \mbox{\texttt{\mdseries\slshape lpgroup}} \cite{MR2864562}. The method reduces the case to certain finitely presented groups by applying
only finitely many endomorphisms to the iterated relations. It then uses coset
enumeration for finitely presented groups to compute an upper bound on the
index of the subgroup. If the coset enumeration for finitely presented groups
terminated, the method attempts to prove that the upper bound is sharp. For
further details we refer to \cite{MR2864562}. 

\subsection{\textcolor{Chapter }{IndexInWholeGroup}}
\logpage{[ 4, 2, 1 ]}\nobreak
\hyperdef{L}{X8014135884DCC53E}{}
{\noindent\textcolor{FuncColor}{$\triangleright$\enspace\texttt{IndexInWholeGroup({\mdseries\slshape H})\index{IndexInWholeGroup@\texttt{IndexInWholeGroup}}
\label{IndexInWholeGroup}
}\hfill{\scriptsize (method)}}\\
\noindent\textcolor{FuncColor}{$\triangleright$\enspace\texttt{FactorCosetAction({\mdseries\slshape G, H})\index{FactorCosetAction@\texttt{FactorCosetAction}}
\label{FactorCosetAction}
}\hfill{\scriptsize (method)}}\\


 The first command attempts to compute the index of \mbox{\texttt{\mdseries\slshape H}} in its parent group. The second one gives the permutation action of \mbox{\texttt{\mdseries\slshape G}} on the right cosets of \mbox{\texttt{\mdseries\slshape H}}. 
\begin{Verbatim}[commandchars=!@|,fontsize=\small,frame=single,label=Example]
  !gapprompt@gap>| !gapinput@G := ExamplesOfLPresentations(1);;|
  !gapprompt@gap>| !gapinput@a := G.1;; b := G.2;; c := G.3;; d := G.4;;|
  !gapprompt@gap>| !gapinput@K := Subgroup( G, [ Comm(a,b), Comm( b, d^a ), Comm( b^a, d )] );;|
  !gapprompt@gap>| !gapinput@IndexInWholeGroup( K );|
  16
  !gapprompt@gap>| !gapinput@FactorCosetAction( G, K );|
  [ a, b, c, d ] -> [ (1,2)(3,6)(4,9)(5,10)(7,11)(8,12)(13,15)(14,16),
    (1,3)(2,6)(4,5)(7,8)(9,10)(11,12)(13,14)(15,16), (1,4)(2,7)(3,5)(6,8)(9,13)(10,14)(11,15)(12,16),
    (1,5)(2,8)(3,4)(6,7)(9,14)(10,13)(11,16)(12,15) ]
\end{Verbatim}
 }

 

\subsection{\textcolor{Chapter }{Index}}
\logpage{[ 4, 2, 2 ]}\nobreak
\hyperdef{L}{X83A0356F839C696F}{}
{\noindent\textcolor{FuncColor}{$\triangleright$\enspace\texttt{Index({\mdseries\slshape H, I})\index{Index@\texttt{Index}}
\label{Index}
}\hfill{\scriptsize (method)}}\\


 attempts to compute the index of \mbox{\texttt{\mdseries\slshape I}} in the subgroup \mbox{\texttt{\mdseries\slshape H}}. The subgroup \mbox{\texttt{\mdseries\slshape I}} must be contained in \mbox{\texttt{\mdseries\slshape H}}. 
\begin{Verbatim}[commandchars=!@|,fontsize=\small,frame=single,label=Example]
  !gapprompt@gap>| !gapinput@G := ExamplesOfLPresentations(1);;|
  !gapprompt@gap>| !gapinput@a := G.1;; b := G.2;; c := G.3;; d := G.4;;|
  !gapprompt@gap>| !gapinput@K := Subgroup( G, [ Comm(a,b), Comm( b, d^a ), Comm( b^a, d )] );;|
  !gapprompt@gap>| !gapinput@KxK := Subgroup( G, [ Comm(b,d^a), Comm(b^a,d), Comm(d^a,c^(a*c)),                         |
  </A> Comm( d^(a*c), c^a), Comm( d, c^(a*c*a) ), Comm( d^(a*c*a), c) ] );;
  !gapprompt@gap>| !gapinput@Index( K, KxK );|
  4
\end{Verbatim}
 }

 

\subsection{\textcolor{Chapter }{CosetTableInWholeGroup}}
\logpage{[ 4, 2, 3 ]}\nobreak
\hyperdef{L}{X846EC8AB7803114D}{}
{\noindent\textcolor{FuncColor}{$\triangleright$\enspace\texttt{CosetTableInWholeGroup({\mdseries\slshape H})\index{CosetTableInWholeGroup@\texttt{CosetTableInWholeGroup}}
\label{CosetTableInWholeGroup}
}\hfill{\scriptsize (method)}}\\


 computes a coset\texttt{\symbol{45}}table for the subgroup \mbox{\texttt{\mdseries\slshape H}} in its parent group. 
\begin{Verbatim}[commandchars=!@|,fontsize=\small,frame=single,label=Example]
  !gapprompt@gap>| !gapinput@CosetTableInWholeGroup( K );|
  [ [ 2, 1, 6, 9, 10, 3, 11, 12, 4, 5, 7, 8, 15, 16, 13, 14 ],
    [ 2, 1, 6, 9, 10, 3, 11, 12, 4, 5, 7, 8, 15, 16, 13, 14 ],
    [ 3, 6, 1, 5, 4, 2, 8, 7, 10, 9, 12, 11, 14, 13, 16, 15 ],
    [ 3, 6, 1, 5, 4, 2, 8, 7, 10, 9, 12, 11, 14, 13, 16, 15 ],
    [ 4, 7, 5, 1, 3, 8, 2, 6, 13, 14, 15, 16, 9, 10, 11, 12 ],
    [ 4, 7, 5, 1, 3, 8, 2, 6, 13, 14, 15, 16, 9, 10, 11, 12 ],
    [ 5, 8, 4, 3, 1, 7, 6, 2, 14, 13, 16, 15, 10, 9, 12, 11 ],
    [ 5, 8, 4, 3, 1, 7, 6, 2, 14, 13, 16, 15, 10, 9, 12, 11 ] ]
\end{Verbatim}
 }

 }

 
\section{\textcolor{Chapter }{Technical details}}\logpage{[ 4, 3, 0 ]}
\hyperdef{L}{X87A9EC0A7DF04931}{}
{


 For performance issues the following global variables can be used to modify
the behaviour of the coset enumeration: 

\subsection{\textcolor{Chapter }{LPRES{\textunderscore}TCSTART}}
\logpage{[ 4, 3, 1 ]}\nobreak
\hyperdef{L}{X823EECA37A8EC3FE}{}
{\noindent\textcolor{FuncColor}{$\triangleright$\enspace\texttt{LPRES{\textunderscore}TCSTART\index{LPRES{\textunderscore}TCSTART@\texttt{LPRES{\textunderscore}TCSTART}}
\label{LPRESuScoreTCSTART}
}\hfill{\scriptsize (global variable)}}\\


 defines the maximal word\texttt{\symbol{45}}length of endomorphisms in the
free monoid which are applied to the iterated relations. }

 

\subsection{\textcolor{Chapter }{LPRES{\textunderscore}CosetEnumerator}}
\logpage{[ 4, 3, 2 ]}\nobreak
\hyperdef{L}{X7C46A9B57BA4CA84}{}
{\noindent\textcolor{FuncColor}{$\triangleright$\enspace\texttt{LPRES{\textunderscore}CosetEnumerator\index{LPRES{\textunderscore}CosetEnumerator@\texttt{LPR}\-\texttt{E}\-\texttt{S{\textunderscore}}\-\texttt{Coset}\-\texttt{Enumerator}}
\label{LPRESuScoreCosetEnumerator}
}\hfill{\scriptsize (global variable)}}\\


 defines the coset enumeration process used for finitely presented groups. It
should be a function which take as input a subgroup \mbox{\texttt{\mdseries\slshape h}} of a finitely presented group and it computes a coset table in the whole
group. The default uses the following method of the \textsf{ACE}\texttt{\symbol{45}}package 
\begin{Verbatim}[commandchars=!@|,fontsize=\small,frame=single,label=Example]
  function ( h )
      local  f, rels, gens;
      f := FreeGeneratorsOfFpGroup( Parent( h ) );
      rels := RelatorsOfFpGroup( Parent( h ) );
      gens := List( GeneratorsOfGroup( h ), UnderlyingElement );
      return ACECosetTable( f, rels, gens : silent := true,
          hard := true,
          max := 10 ^ 8,
          Wo := 10 ^ 8 );
\end{Verbatim}
 If the \textsf{ACE}\texttt{\symbol{45}}package is not available, the library coset enumeration
process is used. }

 }

 }

 
\chapter{\textcolor{Chapter }{Approximating the Schur multiplier}}\logpage{[ 5, 0, 0 ]}
\hyperdef{L}{X7FBE94957D7ECCFC}{}
{
 The algorithm in \cite{MR2678952} approximates the Schur multiplier of an invariantly finitely
L\texttt{\symbol{45}}presented group by the quotients in its
Dwyer\texttt{\symbol{45}}filtration. This is implemented in the \textsf{lpres}\texttt{\symbol{45}}package and the following methods are available: 
\section{\textcolor{Chapter }{Methods}}\logpage{[ 5, 1, 0 ]}
\hyperdef{L}{X8606FDCE878850EF}{}
{
 

\subsection{\textcolor{Chapter }{GeneratingSetOfMultiplier}}
\logpage{[ 5, 1, 1 ]}\nobreak
\hyperdef{L}{X83A5F95E84D3B662}{}
{\noindent\textcolor{FuncColor}{$\triangleright$\enspace\texttt{GeneratingSetOfMultiplier({\mdseries\slshape lpgroup})\index{GeneratingSetOfMultiplier@\texttt{GeneratingSetOfMultiplier}}
\label{GeneratingSetOfMultiplier}
}\hfill{\scriptsize (operation)}}\\


 uses Tietze transformations for computing an equivalent set of relators for \mbox{\texttt{\mdseries\slshape lpgroup}} so that a generating set for its Schur multiplier can be read off easily. }

 

\subsection{\textcolor{Chapter }{FiniteRankSchurMultiplier}}
\logpage{[ 5, 1, 2 ]}\nobreak
\hyperdef{L}{X87A3D6C07D99C79A}{}
{\noindent\textcolor{FuncColor}{$\triangleright$\enspace\texttt{FiniteRankSchurMultiplier({\mdseries\slshape lpgroup, c})\index{FiniteRankSchurMultiplier@\texttt{FiniteRankSchurMultiplier}}
\label{FiniteRankSchurMultiplier}
}\hfill{\scriptsize (operation)}}\\


 computes a finitely generated quotient of the Schur multiplier of \mbox{\texttt{\mdseries\slshape lpgroup}}. The method computes the image of the Schur multiplier of \mbox{\texttt{\mdseries\slshape lpgroup}} in the Schur multiplier of its class\texttt{\symbol{45}}\mbox{\texttt{\mdseries\slshape c}} quotient. }

 

\subsection{\textcolor{Chapter }{EndomorphismsOfFRSchurMultiplier}}
\logpage{[ 5, 1, 3 ]}\nobreak
\hyperdef{L}{X78084374873BDFE1}{}
{\noindent\textcolor{FuncColor}{$\triangleright$\enspace\texttt{EndomorphismsOfFRSchurMultiplier({\mdseries\slshape lpgroup, c})\index{EndomorphismsOfFRSchurMultiplier@\texttt{EndomorphismsOfFRSchurMultiplier}}
\label{EndomorphismsOfFRSchurMultiplier}
}\hfill{\scriptsize (operation)}}\\


 computes a list of endomorphisms of the `FiniteRankSchurMultiplier' of \mbox{\texttt{\mdseries\slshape lpgroup}}. These are the endomorphisms of the invariant
L\texttt{\symbol{45}}presentation induced to `FiniteRankSchurMultiplier'. }

 

\subsection{\textcolor{Chapter }{EpimorphismCoveringGroups}}
\logpage{[ 5, 1, 4 ]}\nobreak
\hyperdef{L}{X7CF92D9880A3687E}{}
{\noindent\textcolor{FuncColor}{$\triangleright$\enspace\texttt{EpimorphismCoveringGroups({\mdseries\slshape lpgroup, d, c})\index{EpimorphismCoveringGroups@\texttt{EpimorphismCoveringGroups}}
\label{EpimorphismCoveringGroups}
}\hfill{\scriptsize (operation)}}\\


 computes an epimorphism of the covering group of the class\texttt{\symbol{45}}\mbox{\texttt{\mdseries\slshape d}} quotient onto the covering group of the class\texttt{\symbol{45}}\mbox{\texttt{\mdseries\slshape c}} quotient. }

 

\subsection{\textcolor{Chapter }{EpimorphismFiniteRankSchurMultiplier}}
\logpage{[ 5, 1, 5 ]}\nobreak
\hyperdef{L}{X86EAE6457CE03B7B}{}
{\noindent\textcolor{FuncColor}{$\triangleright$\enspace\texttt{EpimorphismFiniteRankSchurMultiplier({\mdseries\slshape lpgroup, d, c})\index{EpimorphismFiniteRankSchurMultiplier@\texttt{Epimorphism}\-\texttt{Finite}\-\texttt{Rank}\-\texttt{Schur}\-\texttt{Multiplier}}
\label{EpimorphismFiniteRankSchurMultiplier}
}\hfill{\scriptsize (operation)}}\\


 computes an epimorphism of the $d$\texttt{\symbol{45}}th `FiniteRankSchurMultiplier' of the invariant \mbox{\texttt{\mdseries\slshape lpgroup}} onto the $c$\texttt{\symbol{45}}th `FiniteRankSchurMultiplier'. Its restricts the
epimorphism `EpimorphismCoveringGroups' to the corresponding finite rank
multipliers. }

 

\subsection{\textcolor{Chapter }{ImageInFiniteRankSchurMultiplier}}
\logpage{[ 5, 1, 6 ]}\nobreak
\hyperdef{L}{X87182BC081DCA91E}{}
{\noindent\textcolor{FuncColor}{$\triangleright$\enspace\texttt{ImageInFiniteRankSchurMultiplier({\mdseries\slshape lpgroup, c, elm})\index{ImageInFiniteRankSchurMultiplier@\texttt{ImageInFiniteRankSchurMultiplier}}
\label{ImageInFiniteRankSchurMultiplier}
}\hfill{\scriptsize (function)}}\\


 computes the image of the free group element \mbox{\texttt{\mdseries\slshape elm}} in the \mbox{\texttt{\mdseries\slshape c}}\texttt{\symbol{45}}th `FiniteRankSchurMultiplier'. Note that \mbox{\texttt{\mdseries\slshape elm}} must be a relator contained in the Schur multiplier of \mbox{\texttt{\mdseries\slshape lpgroup}}; otherwise, the function fails in computing the image.

 The following example tackels the Schur multiplier of the Grigorchuk group. 
\begin{Verbatim}[commandchars=!@|,fontsize=\small,frame=single,label=Example]
  !gapprompt@gap>| !gapinput@G := ExamplesOfLPresentations( 1 );;|
  !gapprompt@gap>| !gapinput@gens := GeneratingSetOfMultiplier( G );|
  rec( FixedGens := [ b^-2*c^-2*d^-2*b*c*d*b*c*d ],
    IteratedGens := [ d^-1*a^-1*d^-1*a*d*a^-1*d*a,
        d^-1*a^-1*c^-1*a^-1*c^-1*a^-1*d^-1*a*c*a*c*a*d*a^-1*c^-1*a^-1*c^-1*a^
          -1*d*a*c*a*c*a ],
    BasisGens := [ a^2, b*c*d, b^-2*d^-2*b*c*d*b*c*d, b^-2*c^-2*b*c*d*b*c*d ],
    Endomorphisms := [ [ a, b, c, d ] -> [ a^-1*c*a, d, b, c ] ] )
  !gapprompt@gap>| !gapinput@H := FiniteRankSchurMultiplier( G, 5 );|
  Pcp-group with orders [ 2, 2, 2 ] 
  !gapprompt@gap>| !gapinput@GeneratorsOfGroup( H );|
  [ g15, g17, g16 ]
  !gapprompt@gap>| !gapinput@EndomorphismsOfFRSchurMultiplier( G, 5 );|
  [ [ g15, g16, g17 ] -> [ g15, id, g16 ] ]
  !gapprompt@gap>| !gapinput@Kernel( last[1] );|
  Pcp-group with orders [ 2 ]
  !gapprompt@gap>| !gapinput@GeneratorsOfGroup( last );|
  [ g16 ]
  !gapprompt@gap>| !gapinput@EpimorphismFiniteRankSchurMultipliers( G, 5, 2 );|
  [ g15, g16, g17 ] -> [ g10, id, g13 ]
  !gapprompt@gap>| !gapinput@Range( last ) = FiniteRankSchurMultiplier( G, 2 );|
  true
  !gapprompt@gap>| !gapinput@Kernel( EpimorphismFiniteRankSchurMultipliers( G, 5, 2 ) );|
  Pcp-group with orders [ 2 ]
  !gapprompt@gap>| !gapinput@GeneratorsOfGroup( last );|
  [ g16 ]
  !gapprompt@gap>| !gapinput@Kernel( EpimorphismFiniteRankSchurMultipliers( G, 5, 2 ) ) =|
  </A> Kernel( EndomorphismsOfFRSchurMultiplier( G, 5 )[1] );
  true
  !gapprompt@gap>| !gapinput@ImageInFiniteRankSchurMultiplier( G, 5, gens.FixedGens[1] );|
  g15
  !gapprompt@gap>| !gapinput@ImageInFiniteRankSchurMultiplier(G,5,Image(gens.Endomorphisms[1],|
  </A> gens.IteratedGens[1] ) );
  g16
  !gapprompt@gap>| !gapinput@ImageInFiniteRankSchurMultiplier(G,5,gens.IteratedGens[1] );|
  g17
\end{Verbatim}
 }

 }

 }

 
\chapter{\textcolor{Chapter }{On a parallel nilpotent quotient algorithm}}\logpage{[ 6, 0, 0 ]}
\hyperdef{L}{X7BC16B0082A2B827}{}
{


 We included a parallel version of \textsf{lpres}'s nilpotent quotient algorithm using the \textsf{ParGAP}\texttt{\symbol{45}}package of \textsf{GAP}; see \cite{ParGap}. In this chapter, we outline the basic usage of this parallel part of the \textsf{lpres}\texttt{\symbol{45}}package. For further details on the parallel \textsf{GAP}\texttt{\symbol{45}}sessions we refer to the \textsf{ParGAP}\texttt{\symbol{45}}manual \cite{ParGap}. We note that the \textsf{ParGAP}\texttt{\symbol{45}}package has some bottlenecks in practice. Nevertheless the
significant speed\texttt{\symbol{45}}up of our computations on a
multiple\texttt{\symbol{45}}core system shows that this is a reasonable
extension of the \textsf{lpres}\texttt{\symbol{45}}package. 
\section{\textcolor{Chapter }{Usage}}\logpage{[ 6, 1, 0 ]}
\hyperdef{L}{X86A9B6F87E619FFF}{}
{
 For using the parallel version of the nilpotent quotient algorithm, you will
need to install the \textsf{ParGAP}\texttt{\symbol{45}}package as described in its manual \cite{ParGap}. When using Version 1.1.2 of the \textsf{ParGAP}\texttt{\symbol{45}}package, you will need to apply the following patch to
`pargap/lib/masslave.g' as otherwise the \textsf{ParGAP}\texttt{\symbol{45}}session may crash. On a linux machine you can simply use
`patch {\textless} ../../lpres/gap/pargap/patch' from within the directory
`pargap/lib/'. 
\begin{verbatim}  
  --- masslave.g	2001-11-16 13:17:04.000000000 +0100
  +++ masslave.g	2009-05-06 12:20:19.000000000 +0200
  @@ -467,8 +467,9 @@
     if Length(deltas)>1 then max2 := Maximum(max2, deltas[Length(deltas)-1]); fi;
     max1 := deltas[Length(deltas)];
     pos1 := Position( List(slaveArray, x->realtime-x.time), max1 );
  -  if max1 > slaveTaskTimeFactor and max1 > 30
  -     and slaveTaskTime[pos2].total > 60 then
  +  if max1 > slaveTaskTimeFactor and
  +     max1 > 30 and pos2 <> fail and 
  +     slaveTaskTime[pos2].total > 60 then
       Print("SLAVE ",pos1," SEEMS DEAD!!\n");
     fi;
   end);
\end{verbatim}
 Now, you are ready for creating a \textsf{ParGAP}\texttt{\symbol{45}}session and you can load the \textsf{lpres}\texttt{\symbol{45}}package from within \textsf{ParGAP} using `LoadPackage' as usual. The same methods as described previously are
available. The following example shows the application of the
`NilpotentQuotient'\texttt{\symbol{45}}method to the Grigorchuk group on a
quad\texttt{\symbol{45}}core machine. Note that the significant
speed\texttt{\symbol{45}}up of the nilpotent quotient algorithm is especially
noticeable for large nilpotent quotients. This parallel version of \textsf{lpres} successfully computes some nilpotent quotients which normally took more than a
month to complete. 
\begin{Verbatim}[commandchars=!@|,fontsize=\small,frame=single,label=Example]
  GAP4, Version: 4.4.12 of 17-Dec-2008, i686-pc-linux-gnu-gcc
  GAP4, Version: 4.4.12 of 17-Dec-2008, i686-pc-linux-gnu-gcc
  GAP4, Version: 4.4.12 of 17-Dec-2008, i686-pc-linux-gnu-gcc
  GAP4, Version: 4.4.12 of 17-Dec-2008, i686-pc-linux-gnu-gcc
  GAP4, Version: 4.4.12 of 17-Dec-2008, i686-pc-linux-gnu-gcc
  !gapprompt@gap>| !gapinput@TOPCnumSlaves;|
  4
  !gapprompt@gap>| !gapinput@LoadPackage("LPRES");|
  true
  !gapprompt@gap>| !gapinput@G:=ExamplesOfLPresentations(1);|
  <L-presented group on the generators [ a, b, c, d ]>
  !gapprompt@gap>| !gapinput@SetInfoLevel(InfoLPRES,1);|
  !gapprompt@gap>| !gapinput@NilpotentQuotient(G,2);|
  #I  Class 1: 3 generators with relative orders: [ 2, 2, 2 ]
  #I  Computing a polycyclic presentation for the covering group...
  #I  Checking the consistency relations...
  master -> 1: (AGGLOM_TASK): [ [ -3, 1 ], [ -3, 2 ], [ -2, 1 ], [ 2, -1 ],
    [ 3, -1 ] ]
  master -> 2: (AGGLOM_TASK): [ [ 3, -2 ], [ 1 ], [ 2 ], [ 3 ] ]
  1 -> master: [ [ 0, 0, 0, 0, 0, -2, 0 ], [ 0, 0, 0, 0, 0, 0, -2 ],
    [ 0, 0, 0, 0, -2, 0, 0 ], [ 0, 0, 0, 0, -2, 0, 0 ],
    [ 0, 0, 0, 0, 0, -2, 0 ] ]
  2 -> master: [ [ 0, 0, 0, 0, 0, 0, -2 ], [ 0, 0, 0, 0, 0, 0, 0 ],
    [ 0, 0, 0, 0, 0, 0, 0 ], [ 0, 0, 0, 0, 0, 0, 0 ] ]
  #I  Broadcasting the slaves...
  #I  Inducing the endomorphisms...
  master -> 1:  1
  master -> 2:  2
  master -> 3:  3
  master -> 4:  4
  3 -> master: [ 2, 1 ]
  UPDATE: [ 3, [ 2, 1 ] ]
  1 -> master: [ 2, 1, 8, 1 ]
  UPDATE: [ 1, [ 2, 1, 8, 1 ] ]
  2 -> master: [ 2, 1, 3, 1, 4, 1 ]
  UPDATE: [ 2, [ 2, 1, 3, 1, 4, 1 ] ]
  master -> 1:  5
  master -> 2:  6
  master -> 3:  7
  4 -> master: [ 4, -1, 6, -1, 10, -1 ]
  UPDATE: [ 4, [ 4, -1, 6, -1, 10, -1 ] ]
  1 -> master: [ 6, 1, 8, 2 ]
  UPDATE: [ 5, [ 6, 1, 8, 2 ] ]
  2 -> master: [ 4, 2, 6, 1, 7, 1, 10, 1 ]
  UPDATE: [ 6, [ 4, 2, 6, 1, 7, 1, 10, 1 ] ]
  3 -> master: [ 6, 1 ]
  UPDATE: [ 7, [ 6, 1 ] ]
  master -> 1:  8
  master -> 2:  9
  master -> 3:  10
  1 -> master: [ 10, 1 ]
  UPDATE: [ 8, [ 10, 1 ] ]
  2 -> master: [  ]
  UPDATE: [ 9, [  ] ]
  3 -> master: [ 10, -1 ]
  UPDATE: [ 10, [ 10, -1 ] ]
  #I  Broadcasting the slaves...
  #I  Mapping the relations...
  master -> 1:  1
  master -> 2:  2
  master -> 3:  3
  master -> 4:  4
  1 -> master: [ 0, 0, 0, 0, 1, 0, 0, 0, 0, 0 ]
  2 -> master: [ 0, 0, 0, 2, 0, 1, 1, 0, 0, 1 ]
  3 -> master: [ 0, 0, 0, 0, 0, 1, 0, 0, 0, 0 ]
  4 -> master: [ 0, 0, 0, 0, 0, 0, 1, 0, 0, 0 ]
  master -> 1:  5
  master -> 2:  6
  master -> 3:  7
  1 -> master: [ 0, 0, 0, 1, 0, 1, 1, 0, 0, 1 ]
  2 -> master: [ 0, 0, 0, 0, 0, 0, 0, 0, 2, 0 ]
  3 -> master: [ 0, 0, 0, 0, 0, 0, 0, 4, 6, 4 ]
  #I  Start spinning...
  #I  Extend the quotient system...
  #I  Class 2: 2 generators with relative orders: [ 2, 2 ]
  Pcp-group with orders [ 2, 2, 2, 2, 2 ]
\end{Verbatim}


 Note that the only difference in the parallel version of the \textsf{lpres}\texttt{\symbol{45}}package is a parallel version of the operation
`ExtendQuotientSystem'. This latter operation covers the induction step of the
nilpotent quotient algorithm. }

 }

 \def\bibname{References\logpage{[ "Bib", 0, 0 ]}
\hyperdef{L}{X7A6F98FD85F02BFE}{}
}

\bibliographystyle{alpha}
\bibliography{lpres.bib}

\addcontentsline{toc}{chapter}{References}

\def\indexname{Index\logpage{[ "Ind", 0, 0 ]}
\hyperdef{L}{X83A0356F839C696F}{}
}

\cleardoublepage
\phantomsection
\addcontentsline{toc}{chapter}{Index}


\printindex

\immediate\write\pagenrlog{["Ind", 0, 0], \arabic{page},}
\newpage
\immediate\write\pagenrlog{["End"], \arabic{page}];}
\immediate\closeout\pagenrlog
\end{document}
